\documentclass[12pt,a4paper]{article}
\usepackage[utf8]{inputenc}
\usepackage[german, ngerman]{babel}
\addto\extrasgerman{\languageshorthands{none}}
\usepackage[T1]{fontenc}

\usepackage{geometry}
\usepackage{graphicx}
\usepackage{amsmath}
\usepackage{amssymb}
\usepackage{tikz}
\usepackage{pgfplots}
\usepackage{float}
\usepackage{booktabs}
\usepackage{tabularx}
\usepackage{xcolor}
\usepackage{listings}
\usepackage{microtype}
\usepackage{caption}
\usepackage{subcaption}
\usepackage{hyperref}
\usepackage{acronym}
\usepackage{fancyhdr}
\usepackage{cite}
\usepackage{url}
\usepackage{array}
\usepackage{siunitx}
\usepackage{eurosym}
\usepackage{enumitem}
\usepackage{parskip}
\usepackage{rotating}
\usepackage{enumitem}
\setlist[itemize,1]{label=--}

\hypersetup{colorlinks=true,linkcolor=blue!60!black,urlcolor=blue!60!black,citecolor=blue!60!black}

\usetikzlibrary{shapes.geometric, arrows, arrows.meta, positioning, fit,
  decorations.pathreplacing, backgrounds, calc, patterns, shadows,
  matrix, chains, shapes.multipart, shapes.callouts,
  decorations.markings, automata}
\pgfplotsset{compat=1.18}

\geometry{margin=2.5cm}

\definecolor{etherblue}{RGB}{0,90,170}
\definecolor{appgreen}{RGB}{0,150,80}
\definecolor{warnorange}{RGB}{230,120,0}
\definecolor{lightgray}{RGB}{240,240,240}
\definecolor{darkgray}{RGB}{80,80,80}
\definecolor{accentred}{RGB}{180,30,30}
\definecolor{tabletblue}{RGB}{50,130,200}
\definecolor{cloudblue}{RGB}{100,170,230}

\lstset{
  basicstyle=\ttfamily\small,
  keywordstyle=\color{etherblue}\bfseries,
  commentstyle=\color{appgreen}\itshape,
  stringstyle=\color{accentred},
  breaklines=true,
  frame=single,
  numbers=left,
  numberstyle=\tiny\color{darkgray},
  backgroundcolor=\color{lightgray},
  tabsize=2
}

\hypersetup{
  colorlinks=true,
  linkcolor=etherblue,
  citecolor=appgreen,
  urlcolor=tabletblue,
  pdftitle={EtherCAT Extension: App-Based HMI},
}

\title{
  \textbf{\LARGE EtherCAT Protocol Erweiterung}\\
  \textbf{\Large App-basierte HMI als Ersatz für HMI-Controller}\\[0.4cm]
  \large Mobile und Desktop-Anwendungsschnittstelle\\
   in EtherCAT-Netzwerken\\[0.3cm]
}
\author{Stephan Epp}

\date{\today}

\begin{document}
\maketitle
\thispagestyle{empty}

% =====================================================================

\tableofcontents
\newpage

% =====================================================================
\section{Einleitung}
\label{sec:intro}
% =====================================================================
Diese Arbeit schlägt eine Erweiterung des EtherCAT-Standards vor, die formal eine neue
Klasse von Human-Machine Interfaces (HMI) definiert, die ausschlie\ss{}lich auf mobilen und
Desktop-Applikationen basiert. Traditionelle, band-integrierte HMI-Controller sind teuer,
wartungsintensiv und können die Anzeigequalität moderner Consumer-Hardware bei weitem
nicht erreichen.

Die vorgeschlagene Erweiterung führt einen \emph{App-Based HMI Layer} (ABH) im
EtherCAT-Ökosystem ein, der es Tablets und Smartphones erlaubt, vollständige
HMI-Verantwortung auf Maschinenebene zu übernehmen, während eine erweiterte
Desktop-Variante für Inbetriebnahme und Diagnose durch Applikationsingenieure genutzt
wird. Diese Arbeit beschreibt im Detail die Systemarchitektur, die notwendigen
Protokollerweiterungen, das Sicherheitsmodell mit Authentifizierung und rollenbasierter
Zugriffskontrolle, eine vierstufige Migrationsstrategie sowie ein Referenz-Implementierungsframework.

Die quantitative Kosten-Nutzen-Analyse zeigt eine Einsparung von bis zu 81\,\% über
fünf Jahre bei 20 HMI-Stationen. Der Proposal richtet sich an die EtherCAT Technology
Group (ETG) mit dem Ziel, EHAE als formale Erweiterung des EtherCAT Application Layer
in den ETG.1000-Standard aufzunehmen.

EtherCAT (Ethernet for Control Automation Technology), standardisiert in IEC~61158 und
IEC~61784, ist eines der am weitesten verbreiteten Echtzeit-Industrial-Ethernet-Protokolle
weltweit. Mit seiner segment-basierten Topologie, hochgeschwindigkeits-Datenverarbeitung
und dem verteilten Uhr-Mechanismus (Distributed Clock) ist EtherCAT zur Basis-Infra\-struk\-tur
moderner Maschinenautomation geworden \cite{iec61158}. Zykluszeiten von weniger als
$100\,\si{\micro\second}$ ermöglichen präzise Bewegungssteuerung und synchronisierte
Multi-Achs-Applikationen, die mit klassischen Feldbussystemen nicht realisierbar wären
\cite{etg_spec}.

Die Human-Machine Interface (HMI) --- ein kritisches Teilsystem jeder Maschine --- wurde
jedoch historisch vom Standardisierungsumfang von EtherCAT ausgeschlossen. Dies hat
zur Entstehung einer heterogenen HMI-Landschaft geführt, in der proprietäre Lösungen
verschiedener Hersteller parallel existieren, ohne gemeinsame Schnittstellen oder
Interoperabilität zu gewährleisten.

\subsection{Problemstellung}

Traditionell rüsten Maschinenbauer jede Produktionslinie (sog.\ \emph{Band}) mit einem
dedizierten HMI-Panel aus: einem eingebetteten Touchscreen-System, das eine proprietäre
Runtime ausführt und direkt am Maschinengehäuse montiert ist. Obwohl diese Systeme
ihren Zweck grundsätzlich erfüllen, bringen sie erhebliche Nachteile mit sich:

\begin{itemize}[leftmargin=2em]
  \item \textbf{Hohe Anschaffungskosten:} Industrielle Panel-PCs liegen im Preisbereich
    von \EUR{800} bis über \EUR{5\,000} pro Einheit, abhängig von Schutzklasse,
    Displaygrö\ss{}e und Rechnerleistung.
  \item \textbf{Geringe Anzeigequalität:} Auflösungen zwischen 800$\times$600 und
    1920$\times$1080 bei eingeschränkter Helligkeit (typisch 350--600\,nit) sind der
    Standard.
  \item \textbf{Proprietärer Runtime-Lock-in:} Jeder Hersteller nutzt eine eigene
    Entwicklungsumgebung und Runtime (z.\,B.\ Siemens WinCC, Beckhoff TwinCAT HMI,
    Wein\-tek EasyBuilder), was Portierbarkeit und Herstellerunabhängigkeit stark einschränkt.
  \item \textbf{Wartungsaufwand:} Ersatzteilbeschaffung und Software-Kompatibilität bei
    veralteten Panels stellen langfristig erhebliche Herausforderungen dar.
  \item \textbf{Fehlende Konnektivität:} Integrierte Remote-Monitoring- und
    Cloud-Anbindungen sind nur mit proprietären Zusatzlösungen realisierbar.
  \item \textbf{Langer Update-Zyklus:} Firmware- und Software-Updates erfordern oft
    manuelle Eingriffe vor Ort und sind nicht zentral ausrollbar.
\end{itemize}

Gleichzeitig übersteigen Consumer-Tablets die HMI-Anzeigequalität bei einem Bruchteil
der Kosten. Ein modernes iPad Pro erreicht Auflösungen von 2752$\times$2064 Pixeln,
Bildwiederholraten von 120\,Hz sowie eine Helligkeit von bis zu 1600\,nit --- Werte, die
industrielle Panels nicht annähernd bieten.

\subsection{Ziel dieser Arbeit}

Diese Arbeit schlägt die \textbf{EtherCAT-HMI-App Extension (EHAE)} vor und definiert,
wie app-basierte HMI-Clients mit EtherCAT-Mastern integriert werden können, um
band-integrierte HMI-Controller vollständig zu ersetzen. Die zentralen Beiträge dieser
Arbeit sind:

\begin{enumerate}[leftmargin=2em]
  \item Definition eines neuen Software-Layers (\emph{HMI Application Gateway, HAG})
    im EtherCAT-Master-Stack.
  \item Spezifikation eines standardisierten Kommunikationsprotokolls (\emph{HAG
    Communication Protocol, HCP}) auf Basis von WebSockets.
  \item Formalisierung eines Dual-Role-App-Modells mit klar definierten
    Berechtigungsstufen.
  \item Beschreibung eines auf OAuth~2.0, TLS~1.3 und RBAC aufbauenden
    Sicherheitsmodells.
  \item Quantitativer Nachweis der Kosteneinsparungen gegenüber traditionellen Ansätzen.
\end{enumerate}

\subsection{Kostenvergleich}

Abbildung~\ref{fig:cost_comparison} illustriert den 5-Jahres-Gesamtbetriebskostenvergleich
(Total Cost of Ownership, TCO) zwischen einem traditionellen HMI-Panel und einer
app-basierten Lösung.

\begin{figure}[H]
\centering
\begin{tikzpicture}
\begin{axis}[
  ybar, bar width=1.2cm, width=13cm, height=7cm,
  ylabel={Kosten (EUR)}, xlabel={Kostenkategorie},
  xtick=data,
  xticklabels={Hardware, Lizenzierung, Wartung, Ersatz, Gesamt},
  xticklabel style={rotate=15, anchor=east},
  legend style={at={(0.99,0.99)}, anchor=north east},
  ymin=0, ymax=22000,
  ymajorgrids=true, grid style=dashed,
  title={5-Jahres-TCO-Vergleich pro HMI-Station},
  title style={font=\bfseries},
]
\addplot[fill=etherblue!80, draw=etherblue] coordinates {
  (1,3500)(2,2200)(3,3000)(4,2000)(5,10700)};
\addplot[fill=appgreen!80, draw=appgreen] coordinates {
  (1,600)(2,200)(3,300)(4,200)(5,1300)};
\legend{Traditionelles HMI-Panel, App-Based HMI (Tablet)}
\end{axis}
\end{tikzpicture}
\caption{5-Jahres-TCO-Vergleich: traditionelles HMI-Panel vs.\ app-basiertes HMI pro
  Station. Das Balkendiagramm zeigt fünf Kostenkategorien (Hardware, Lizenzierung,
  Wartung, Ersatz und Gesamtkosten) jeweils nebeneinander für beide Varianten. Beim
  traditionellen HMI-Panel dominieren die Hardware-Kosten mit \EUR{3\,500} sowie
  Wartungskosten von \EUR{3\,000}. Die app-basierte Lösung reduziert alle
  Einzelpositionen drastisch: Hardware auf \EUR{600}, Wartung auf \EUR{300}. In der
  Gesamtbetrachtung über fünf Jahre ergibt sich \EUR{10\,700} gegenüber nur \EUR{1\,300}
  --- eine Einsparung von \textbf{87\,\%} pro Station.}
\label{fig:cost_comparison}
\end{figure}

Die Differenz zwischen den Varianten ist auf mehrere Faktoren zurückzuführen: Industrielle
Panels erfordern Wartungsverträge, herstellerspezifische Schulungen und kostspielige
Ersatzteile. App-basierte Systeme auf Consumer-Hardware profitieren von der Skalierung
des Massenmarktes, standardisierten App-Store-Distributionen sowie kostenlosen oder
preiswerten Entwicklungsframeworks wie Flutter oder .NET MAUI.

% =====================================================================
\section{Stand der Technik}
\label{sec:sota}
% =====================================================================

\subsection{EtherCAT-Protokollübersicht}

EtherCAT operiert auf einem Master-Slave-Prinzip \cite{etg_spec}. Ein einzelner Master
kommuniziert mit mehreren Slave-Geräten in einer Ring- oder Linientopologie über
Standard-Ethernet-Frames und erreicht dabei Zykluszeiten von weniger als
$100\,\si{\micro\second}$. Das Protokoll verwendet einen innovativen
\emph{Processing-on-the-fly}-Ansatz: Jeder Slave liest die für ihn relevanten Daten aus
dem durchlaufenden Frame heraus und schreibt seine Prozessdaten hinein, während der
Frame den Knoten passiert --- ohne Pufferung und ohne zusätzliche Latenz.

Die EtherCAT-Protokollarchitektur besteht aus drei wesentlichen Schichten:

\begin{enumerate}[leftmargin=2em]
  \item \textbf{Physical Layer (IEC~61158-2):} Definiert die elektrischen und mechanischen
    Eigenschaften der Übertragung, typischerweise 100BASE-TX über Cat-5-Kabel.
  \item \textbf{Data Link Layer:} Beinhaltet das Field Memory Management Unit (FMMU) und
    den SyncManager. Das FMMU ermöglicht das flexible Mapping von Prozessdaten in den
    EtherCAT-Frame, während der SyncManager die synchronisierte Kommunikation zwischen
    Master und Slave verwaltet.
  \item \textbf{Application Layer:} Definiert höhere Protokolle wie CAN over EtherCAT
    (CoE), File over EtherCAT (FoE), Ethernet over EtherCAT (EoE) und ADS over EtherCAT
    (AoE) \cite{iec61784}.
\end{enumerate}

\subsection{Standardtopologie und HMI-Anbindung}

Abbildung~\ref{fig:ethercat_topology} zeigt die typische EtherCAT-Netzwerktopologie mit
traditionellem HMI-Panel.

\begin{figure}[H]
\centering
\begin{tikzpicture}[
  master/.style={rectangle, rounded corners=4pt, draw=etherblue, fill=etherblue!15,
    minimum width=3cm, minimum height=1cm, font=\small\bfseries, text=etherblue, align=center},
  slave/.style={rectangle, rounded corners=3pt, draw=darkgray, fill=lightgray,
    minimum width=2cm, minimum height=0.8cm, font=\footnotesize},
  hmi/.style={rectangle, rounded corners=3pt, draw=accentred, fill=accentred!10,
    minimum width=2cm, minimum height=0.8cm, font=\footnotesize, text=accentred, align=center},
  arr/.style={-Stealth, thick, color=etherblue},
]
\node[master] (master) at (0,0)   {EtherCAT\\Master};
\node[slave]  (s1)     at (4,0)   {Drive 1};
\node[slave]  (s2)     at (7,0)   {Drive 2};
\node[slave]  (s3)     at (10,0)  {I/O Module};
\node[slave]  (s4)     at (7,-2.5){Sensor Hub};
\node[hmi]    (hmi)    at (0,-2.5){HMI Panel\\(traditional)};
\draw[arr] (master.east) -- (s1.west) node[midway, above, font=\tiny] {Ethernet};
\draw[arr] (s1.east) -- (s2.west);
\draw[arr] (s2.east) -- (s3.west);
\draw[arr, dashed] (s3.south) -- ++(0,-0.6) -| (s4.east);
\draw[arr, dashed, color=accentred] (master.south) -- (hmi.north)
  node[midway, right, font=\tiny, color=accentred, align=center] {Modbus/OPC-UA\\(non-standard)};
\node[font=\tiny, color=accentred, align=center] at (-1.2,-1.5) {Proprietary\\Interface};
\end{tikzpicture}
\caption{Standard-EtherCAT-Topologie mit traditionell angebundenem HMI-Panel. Das
  Diagramm zeigt den EtherCAT-Master (blau, links oben) als zentralen Kommunikationsknoten,
  der über eine lineare Ethernet-Kette mit drei Slave-Geräten verbunden ist: Drive~1,
  Drive~2 und einem I/O-Modul (alle grau). Ein vierter Slave, der Sensor Hub, ist mit
  gestrichelter Linie als Abzweigung von Drive~2 dargestellt. Das traditionelle HMI-Panel
  (rot markiert, unten links) ist \emph{nicht} Teil des EtherCAT-Rings, sondern über eine
  gestrichelte, rot markierte Verbindung mit nicht-standardisierten Protokollen (Modbus oder
  OPC-UA) am Master angebunden. Dies illustriert das Kernproblem: HMI ist kein nativer
  Bestandteil des EtherCAT-Standards, sondern eine proprietäre Ergänzung.}
\label{fig:ethercat_topology}
\end{figure}

Diese Entkopplung hat weitreichende Konsequenzen: Da das HMI au\ss{}erhalb des
EtherCAT-Standards liegt, existiert keine einheitliche Schnittstelle. Jeder Hersteller
implementiert die Verbindung zwischen Master und HMI individuell, was zu
Kompatibilitätsproblemen, proprietären Protokollen und erhöhtem Integrationsaufwand führt.

\subsection{Consumer-Hardware-Fähigkeiten}

Der rasante Fortschritt im Consumer-Elektronik-Markt hat eine technologische Lücke
gegenüber industriellen HMI-Panels entstehen lassen, die sich Jahr für Jahr vergrö\ss{}ert.
Tabelle~\ref{tab:hw_comparison} fasst die wichtigsten Parameter zusammen.

\begin{table}[H]
\centering
\caption{Displayvergleich: industrielles HMI-Panel vs.\ Consumer-Tablet (2024).}
\label{tab:hw_comparison}
\begin{tabularx}{\textwidth}{lXX}
\toprule
\textbf{Parameter} & \textbf{Industrielles HMI-Panel} & \textbf{Consumer-Tablet (iPad Pro /
Galaxy Tab S9)} \\
\midrule
Auflösung          & 1024$\times$768 -- 1920$\times$1080  & 2560$\times$1600 --
2752$\times$2064 \\
Bildwiederholrate  & 60\,Hz                               & 60--120\,Hz (ProMotion) \\
Helligkeit         & 350--600\,nit                        & 600--1600\,nit \\
Touch-Präzision    & $\pm$2--5\,mm                        & $\pm$0{,}5\,mm \\
Preis (EUR)        & 800--5\,000                          & 350--1\,300 \\
Prozessor          & ARM Cortex-A, 1--2\,GHz              & Apple M4 / Snapdragon 8 Gen 2 \\
RAM                & 512\,MB -- 2\,GB                     & 8--16\,GB \\
Akkulaufzeit       & entfällt (Netzbetrieb)               & 10--12\,Stunden \\
\bottomrule
\end{tabularx}
\end{table}

Neben den Displayeigenschaften sind Consumer-Tablets in weiteren Dimensionen deutlich
überlegen: ihre Prozessoren ermöglichen flüssige Echtzeitvisualisierungen, Trend-Displays
und sogar lokale Datenanalyse. Die Integration von Wi-Fi~6, Bluetooth und Mobilfunk bietet
Konnektivitätsoptionen, die industrielle Panels nicht besitzen.

\subsection{Bestehende Ansätze und deren Grenzen}

Bereits heute gibt es Versuche, Consumer-Geräte als HMI einzusetzen. Diese basieren
typischerweise auf OPC-UA-Konnektivität oder proprietären Herstellerlösungen wie
Siemens MindSphere oder Beckhoff TF2000. Sie weisen jedoch folgende Schwächen auf:

\begin{itemize}[leftmargin=2em]
  \item Keine formale Standardisierung im EtherCAT-Kontext.
  \item Fehlende Definition von Sicherheitsrollen und Zugriffsrechten.
  \item Keine spezifizierte Latenzanforderung für die HMI-App-Konnektivität.
  \item Vendor-Lock-in durch proprietäre Protokolle.
\end{itemize}

EHAE schlie\ss{}t diese Lücke durch eine vollständige, herstellerunabhängige Spezifikation.

% =====================================================================
\section{EHAE-Architektur}
\label{sec:ehae}
% =====================================================================

\subsection{Systemübersicht}

Die EHAE definiert einen neuen Software-Layer --- den \emph{HMI Application Gateway
(HAG)} --- innerhalb des EtherCAT-Master-Stacks. Der HAG exponiert eine abgesicherte,
WebSocket-basierte API, mit der HMI-App-Clients über Wi-Fi oder kabelgebundenes Ethernet
verbinden können.

Abbildung~\ref{fig:ehae_architecture} zeigt das vollständige Schichtenmodell der EHAE.

\begin{figure}[H]
\centering
\begin{tikzpicture}[
  layer/.style={rectangle, rounded corners=5pt, draw=#1, fill=#1!12,
    minimum width=12cm, minimum height=1cm, font=\small\bfseries, text=#1},
  box/.style={rectangle, rounded corners=3pt, draw=#1, fill=#1!18,
    font=\footnotesize, text=#1, minimum width=2.5cm, minimum height=0.7cm, align=center},
  arr/.style={-Stealth, thick},
]
\node[layer=darkgray]   (phy) at (0,0)   {EtherCAT Physical Layer (IEC 61158-2)};
\node[layer=etherblue]  (dll) at (0,1.4) {EtherCAT Data Link Layer --- FMMU / SyncManager};
\node[layer=appgreen]   (eal) at (0,2.8) {EtherCAT Application Layer (CoE, FoE, EoE, AoE)};
\node[layer=warnorange] (hag) at (0,4.2) {\textbf{[NEU]} HMI Application Gateway (HAG)
  --- EHAE-Erweiterung};
\node[box=tabletblue] (opapp)  at (-4.5,6) {Operator App\\(Tablet/Smartphone)};
\node[box=tabletblue] (engapp) at (-1.5,6) {Engineer App\\(Desktop/Tablet)};
\node[box=cloudblue]  (remapp) at ( 1.5,6) {Remote App\\(Cloud Monitor)};
\node[box=darkgray]   (third)  at ( 4.5,6) {3rd Party\\SCADA};
\foreach \n in {opapp, engapp, remapp, third}
  \draw[arr, color=tabletblue] (\n.south) -- ++(0,-0.4) |- (hag.north);
\draw[arr, color=etherblue, <->] (phy.north) -- (dll.south);
\draw[arr, color=etherblue, <->] (dll.north) -- (eal.south);
\draw[arr, color=appgreen,  <->] (eal.north) -- (hag.south);
\draw[decorate, decoration={brace,amplitude=8pt}, color=accentred]
  (6.9,3.5) -- (6.9,5.0) node[midway, right=6pt, font=\tiny, color=accentred, align=left]
  {TLS 1.3\\OAuth 2.0\\RBAC};
\end{tikzpicture}
\caption{EHAE-Architekturdiagramm: Schichtenmodell der erweiterten EtherCAT-Protokollstack.
  Das Diagramm ist von unten nach oben aufgebaut und zeigt vier aufeinander aufbauende
  Schichten. Ganz unten (grau) liegt der \emph{EtherCAT Physical Layer} gemä\ss{} IEC~61158-2.
  Darüber (blau) befindet sich der \emph{Data Link Layer} mit FMMU und SyncManager.
  Die dritte Schicht (grün) ist der \emph{Application Layer} mit CoE, FoE, EoE und AoE.
  Die vierte, orange hervorgehobene Schicht ist die neue EHAE-Erweiterung: der
  \emph{HMI Application Gateway (HAG)}. Von dieser Schicht aus verbinden sich vier Typen
  von App-Clients über WebSocket-Verbindungen (Pfeile von oben): die Operator~App
  (Tablet/Smartphone), die Engineer~App (Desktop/Tablet), die Remote~App (Cloud-Monitor)
  sowie Drittanbieter-SCADA-Systeme. Rechts neben dem HAG zeigt eine geschweifte Klammer
  die eingesetzten Sicherheitsmechanismen: TLS~1.3, OAuth~2.0 und RBAC.}
\label{fig:ehae_architecture}
\end{figure}

Die zentrale Design-Entscheidung des EHAE-Architekturmodells besteht darin, den HAG
\emph{oberhalb} des bestehenden Application Layers zu positionieren. Dies hat mehrere
Vorteile:

\begin{itemize}[leftmargin=2em]
  \item \textbf{Keine Änderungen an Slave-Firmware:} Alle Erweiterungen laufen
    ausschlie\ss{}lich im Master-Stack.
  \item \textbf{Rückwärtskompatibilität:} Bestehende EtherCAT-Netzwerke können ohne
    Modifikation der Feldgeräte auf EHAE migrieren.
  \item \textbf{Herstellerfreiheit:} Jeder EtherCAT-Master-Stack-Anbieter kann den HAG
    unabhängig implementieren.
\end{itemize}

\subsection{Interne Struktur des HMI Application Gateway}

Abbildung~\ref{fig:hag_internal} zeigt die interne Modulstruktur des HAG.

\begin{sidewaysfigure}
\centering
\begin{tikzpicture}[
  mod/.style={rectangle, rounded corners=4pt, draw=etherblue, fill=etherblue!10,
    minimum width=3cm, minimum height=0.9cm, font=\footnotesize\bfseries, text=etherblue,
    align=center},
  ext/.style={rectangle, rounded corners=4pt, draw=appgreen, fill=appgreen!10,
    minimum width=2.8cm, minimum height=0.8cm, font=\footnotesize, align=center},
  store/.style={cylinder, shape border rotate=90, draw=darkgray, fill=lightgray,
    minimum width=1.8cm, minimum height=0.8cm, font=\footnotesize},
  arr/.style={-Stealth, thick, color=etherblue},
]
\node[mod]   (pdm)   at (0,0)     {PDO\\Mapper};
\node[mod]   (wsrv)  at (4,0)     {WebSocket\\Server};
\node[mod]   (auth)  at (4,2)     {Auth\\Broker};
\node[mod]   (rbac)  at (4,-2)    {RBAC\\Engine};
\node[mod]   (ebus)  at (0,-2)    {Event\\Bus};
\node[mod]   (alm)   at (-4,-2)   {Alarm\\Manager};
\node[mod]   (cfg)   at (-4,0)    {Config\\Loader};
\node[store] (vardb) at (0,2)     {Var DB};
\node[ext]   (eal)   at (-7.5,0)  {EtherCAT\\App Layer};
\node[ext]   (apps)  at (7.5,0)   {App\\Clients};
\draw[arr] (eal)   -- (cfg);
\draw[arr] (cfg)   -- (pdm);
\draw[arr] (pdm)   -- (vardb);
\draw[arr] (vardb) -- (wsrv);
\draw[arr] (wsrv)  -- (apps);
\draw[arr, <->] (wsrv) -- (auth);
\draw[arr, <->] (wsrv) -- (rbac);
\draw[arr] (alm)   -- (ebus);
\draw[arr] (ebus)  -- (wsrv);
\draw[arr] (pdm)   -- (ebus);
\end{tikzpicture}
\caption{Interne Modulstruktur des HAG. Das Blockdiagramm zeigt sieben funktionale Module
  (blau) sowie eine zentrale Datenbank (grau, Zylinder-Symbol) und zwei externe
  Schnittstellen (grün). Von links kommend empfängt der \emph{Config Loader} die
  Prozessdaten-Mappings aus dem EtherCAT Application Layer. Der \emph{PDO Mapper}
  übersetzt rohe EtherCAT-Prozessdaten in benannte HMI-Variablen und speichert diese in
  der \emph{Variable Database (Var DB)}. Der \emph{WebSocket Server} in der Mitte ist
  das Verbindungsstück zu den App-Clients (rechts). Er kommuniziert bidirektional mit dem
  \emph{Auth Broker} (Token-Validierung) und der \emph{RBAC Engine}
  (Zugriffsrechtsprüfung). Der \emph{Event Bus} sammelt Ereignisse aus dem PDO Mapper
  und dem \emph{Alarm Manager} und leitet sie an den WebSocket Server weiter, der sie als
  Push-Nachrichten an die Clients ausliefert.}
\label{fig:hag_internal}
\end{sidewaysfigure}

Jedes Modul des HAG hat eine klar definierte Verantwortlichkeit:

\begin{description}[leftmargin=2em]
  \item[Config Loader] Liest und parst die HMI Variable Map (HVM, siehe
    Abschnitt~\ref{sec:protocol}) und konfiguriert daraus den PDO Mapper zur Laufzeit.
  \item[PDO Mapper] Führt das zyklische Mapping von EtherCAT-Prozessdatenobjekten
    (PDOs) auf benannte HMI-Variablen durch. Erkennt Werteänderungen und generiert
    entsprechende Events.
  \item[Variable Database] In-Memory-Datenbank, die den aktuellen Zustand aller
    HMI-Varia\-blen hält. Dient als konsistenter Snapshot-Speicher für neue
    WebSocket-Verbindungen.
  \item[WebSocket Server] Verwaltet alle aktiven Client-Verbindungen, verarbeitet
    eingehende Nachrichten und verteilt ausgehende Updates asynchron.
  \item[Auth Broker] Validiert JWT-Tokens bei Verbindungsaufbau via
    Token-Introspection-Endpunkt des Identity Providers.
  \item[RBAC Engine] Prüft vor jeder schreibenden Operation, ob die Rolle des Clients die
    erforderliche Berechtigung besitzt.
  \item[Event Bus] Internes Publish-Subscribe-System für lose Kopplung zwischen Modulen.
  \item[Alarm Manager] Verwaltet den Alarm-Lifecycle (raise, acknowledge, clear) und
    persistiert die Alarm-Historie.
\end{description}

% =====================================================================
\section{Dual-Role-App-Modell}
\label{sec:dualrole}
% =====================================================================

EHAE definiert formal zwei App-Rollen, die auf demselben Anwendungsframework laufen,
aber unterschiedliche Capability-Sets besitzen. Dieses Modell erlaubt es, eine einzige
Codebase für beide Rollen zu nutzen und die Fähigkeiten dynamisch auf Basis des
OAuth-Tokens zur Laufzeit zu aktivieren.

\begin{figure}[H]
\centering
\begin{tikzpicture}[
  phone/.style={rounded corners=8pt, draw=tabletblue, fill=tabletblue!10,
    minimum width=2.2cm, minimum height=3.8cm},
  tablet/.style={rounded corners=6pt, draw=appgreen, fill=appgreen!10,
    minimum width=4cm, minimum height=3cm},
  desktop/.style={rounded corners=4pt, draw=darkgray, fill=lightgray,
    minimum width=6cm, minimum height=3.5cm},
  lbl/.style={font=\footnotesize\bfseries, align=center},
  feat/.style={font=\tiny, align=left},
]
\node[phone] (ph) at (0,0) {};
\node[lbl, text=tabletblue] at (0,1.5)  {Operator\\App};
\node[feat] at (0, 0.4) {\checkmark~Maschine EIN/AUS};
\node[feat] at (0,-0.1) {\checkmark~Statusanzeige};
\node[feat] at (0,-0.6) {\checkmark~Alarm ACK};
\node[font=\tiny, color=accentred] at (0,-1.1) {\texttimes~Parameteränderung};
\node[tablet] (tb) at (5.5,0) {};
\node[lbl, text=appgreen] at (5.5,1.0)  {Operator App\\(Tablet)};
\node[feat] at (5.5, 0.3) {\checkmark~Alle Smartphone-Features};
\node[feat] at (5.5,-0.2) {\checkmark~Trend-Anzeige};
\node[feat] at (5.5,-0.7) {\checkmark~Rezept-Auswahl};
\node[desktop] (dt) at (11.5,0) {};
\node[lbl, text=darkgray] at (11.5,1.5) {Engineer App (Desktop/Tablet)};
\node[feat] at (11.5, 0.8) {\checkmark~Alle Operator-Features};
\node[feat] at (11.5, 0.3) {\checkmark~SDO-Parametrierung};
\node[feat] at (11.5,-0.2) {\checkmark~Antriebsinbetriebnahme};
\node[feat] at (11.5,-0.7) {\checkmark~Firmware-Update (FoE)};
\node[feat] at (11.5,-1.2) {\checkmark~Oszilloskop-Ansicht};
\end{tikzpicture}
\caption{Dual-Role-App-Modell: Gegenüberstellung von Operator~App und Engineer~App.
  Das Diagramm zeigt drei Geräteklassen von links nach rechts, deren Capability-Sets
  streng hierarchisch aufgebaut sind. Links (blau) steht die \emph{Operator App} auf
  einem Smartphone: sie erlaubt grundlegende Steuerungsoperationen (EIN/AUS), die
  Statusanzeige und das Quittieren von Alarmen; Parameteränderungen sind explizit
  gesperrt (\texttimes). In der Mitte (grün) die \emph{Operator App} auf einem Tablet
  mit erweiterter Anzeigefläche: alle Smartphone-Features plus Trendanzeige und
  Rezeptauswahl sind verfügbar. Rechts (grau) die \emph{Engineer App} auf Desktop
  oder Tablet: das vollständige Capability-Set einschlie\ss{}lich SDO-Parametrierung,
  Antriebsinbetriebnahme, Firmware-Updates über FoE und Echtzeitoszilloskop.}
\label{fig:dual_role}
\end{figure}

\subsection{Engineer App: Topologie-Ansicht}

Eine der leistungsstärksten Funktionen der Engineer App ist die Live-Topologie-Visuali\-sierung.
Sie zeigt das gesamte EtherCAT-Netzwerk in Echtzeit, inklusive Gerätezustände, Diagnose-
informationen und Netzwerkstatistiken.

\begin{figure}[H]
\centering
\begin{tikzpicture}[
  master/.style={rectangle, rounded corners=5pt, draw=etherblue, fill=etherblue!20,
    minimum width=2.5cm, minimum height=1cm, font=\small\bfseries, text=etherblue},
  slave/.style={rectangle, rounded corners=3pt, draw=darkgray, fill=lightgray,
    minimum width=2cm, minimum height=0.8cm, font=\scriptsize, align=center},
  ok/.style={circle, draw=appgreen, fill=appgreen!20, minimum size=0.4cm},
  err/.style={circle, draw=accentred, fill=accentred!20, minimum size=0.4cm},
  arr/.style={-Stealth, thick, color=etherblue},
]
\node[master] (m) at (0,0) {Master};
\node[slave] (d1) at (3.5, 1.2) {Drive 1\\Achse X};
\node[slave] (d2) at (3.5,-1.2) {Drive 2\\Achse Y};
\node[slave] (d3) at (7,   1.2) {Drive 3\\Achse Z};
\node[slave] (io1) at (7,  -1.2){I/O Coupler\\EK1100};
\node[slave] (io2) at (10.5,-1.2){I/O Modul\\EL2008};
\node[slave] (enc) at (10.5, 1.2){Encoder\\EL5151};
\node[ok]  at (3.5, 1.85) {}; \node[font=\tiny, color=appgreen]  at (3.5,2.2)  {OP};
\node[ok]  at (3.5,-0.55) {}; \node[font=\tiny, color=appgreen]  at (3.5,-0.2) {OP};
\node[err] at (7,   1.85) {}; \node[font=\tiny, color=accentred] at (7,2.2)    {ERR};
\node[ok]  at (7,  -0.55) {};
\node[ok]  at (10.5,-0.55) {};
\node[ok]  at (10.5, 1.85) {};
\draw[arr] (m.east) -- ++(0.5,0) |- (d1.west);
\draw[arr] (m.east) -- ++(0.5,0) |- (d2.west);
\draw[arr] (d1.east) -- (d3.west);
\draw[arr] (d2.east) -- (io1.west);
\draw[arr] (io1.east) -- (io2.west);
\draw[arr] (d3.south) -- ++(0,-0.4) -- ++(3.5,0) |- (enc.west);
\node[font=\tiny, color=etherblue, align=center] at (0,-2.5)
  {Netzwerkzyklus: 250\,\si{\micro\second} | Working Counter: 12/12 | Frame Loss: 0};
\end{tikzpicture}
\caption{Engineer-App-Topologieansicht: Live-EtherCAT-Netzwerkkarte mit Gerätezuständen.
  Das Diagramm zeigt den EtherCAT-Master (blau) links, von dem zwei parallele Stränge
  abzweigen. Der obere Strang führt über Drive~1 (Achse~X, grün $=$ Betrieb) zu Drive~3
  (Achse~Z, \textbf{rot $=$ Fehler}), der untere über Drive~2 (Achse~Y) zu I/O-Coupler
  EK1100 und I/O-Modul EL2008. Der Encoder EL5151 ist als Abzweigung von Drive~3
  dargestellt. Über jedem Slave-Knoten ist ein farbiger Kreis mit Statuskürzel
  dargestellt: \textbf{grün OP} (Operational) für alle funktionierenden Geräte,
  \textbf{rot ERR} für Drive~3~(Achse~Z), das einen Fehler meldet. Die
  Statuszeile am unteren Rand zeigt Netzwerkparameter: Zykluszeit 250\,µs, Working
  Counter 12/12 (alle Slaves antworten), Frame Loss~0.}
\label{fig:eng_topo_view}
\end{figure}

Der Working Counter (WC) ist ein kritischer Diagnoseparameter in EtherCAT-Netzwerken.
Er zählt, wie viele Slaves korrekt auf den umlaufenden Frame geantwortet haben. Ein WC
von 12/12 bedeutet, dass alle erwarteten 12 Datagramm-Antworten eingetroffen sind.
Abweichungen deuten auf Kommunikationsfehler oder ausgefallene Slaves hin. Die
Engineer App visualisiert diesen Parameter dauerhaft und schlägt Alarm bei
Unstimmigkeiten.

\subsection{Engineer App: Oszilloskop-Ansicht}

Die Engineer App bietet eine vollständige Oszilloskop-Funktion, die Prozessvariablen
in Echtzeit mit einer konfigurierbaren Abtastrate aufzeichnet. Typische Anwendungsfälle
sind:

\begin{itemize}[leftmargin=2em]
  \item Analyse von Regelkreisverhalten (Überschwingen, Einregelzeit, Stationäre
    Abweichung).
  \item Diagnose von mechanischen Resonanzen durch FFT-Analyse.
  \item Inbetriebnahme von Lageregelungen (Position, Geschwindigkeit, Strom).
  \item Aufzeichnung von Anlaufkurven nach Parameteränderungen.
\end{itemize}

% =====================================================================
\section{Protokoll-Erweiterungsspezifikation}
\label{sec:protocol}
% =====================================================================

\subsection{HAG Communication Protocol (HCP)}

Die EHAE definiert das \emph{HAG Communication Protocol (HCP)}, das über WebSocket
(RFC~6455) \cite{rfc6455} mit JSON- oder MessagePack-Serialisierung arbeitet. Die Wahl
von WebSocket über alternatives HTTP/REST ist bewusst getroffen: WebSocket unterstützt
Full-Duplex-Kommunikation und Server-seitige Push-Events, was für eine reaktive HMI mit
geringer Latenz unverzichtbar ist.

Tabelle~\ref{tab:hcp_messages} listet alle definierten HCP-Nachrichtentypen auf.

\begin{table}[H]
\centering
\caption{HCP-Nachrichtentypen.}
\label{tab:hcp_messages}
\begin{tabularx}{\textwidth}{llX}
\toprule
\textbf{Nachricht}          & \textbf{Richtung}     & \textbf{Beschreibung} \\
\midrule
\texttt{SUBSCRIBE}          & Client $\to$ HAG      & Abonnement benannter HMI-Variablen \\
\texttt{DATA\_UPDATE}       & HAG $\to$ Client      & Push aktualisierter Variablenwerte \\
\texttt{WRITE\_VAR}         & Client $\to$ HAG      & Variable schreiben (rollengesperrt) \\
\texttt{ALARM\_EVENT}       & HAG $\to$ Client      & Alarm ausgelöst oder quittiert \\
\texttt{ALARM\_ACK}         & Client $\to$ HAG      & Operator quittiert Alarm \\
\texttt{SDO\_READ/WRITE}    & Client $\to$ HAG      & SDO-Zugriff (nur Engineer-Rolle) \\
\texttt{TOPO\_REQUEST}      & Client $\to$ HAG      & Topologie-Snapshot anfordern \\
\texttt{HEARTBEAT}          & Bidirektional         & Keep-Alive mit Zeitstempel \\
\texttt{SCOPE\_SUBSCRIBE}   & Client $\to$ HAG      & Oszilloskop-Stream starten \\
\texttt{SCOPE\_DATA}        & HAG $\to$ Client      & Hochfrequente Messdaten (bis 1\,kHz) \\
\texttt{FOE\_UPLOAD}        & Client $\to$ HAG      & Firmware-Upload initiieren (FoE) \\
\texttt{FOE\_STATUS}        & HAG $\to$ Client      & Upload-Fortschritt und Status \\
\bottomrule
\end{tabularx}
\end{table}

\subsubsection{Nachrichtenformat}

Jede HCP-Nachricht folgt einem einheitlichen JSON-Schema:

\begin{lstlisting}[language={}, caption={HCP-Nachrichtenformat (Beispiel: DATA\_UPDATE).}]
{
  "type": "DATA_UPDATE",
  "seq": 1042,
  "ts": 1712345678.342,
  "variables": [
    {
      "name": "axis_x.actual_velocity",
      "value": 1247.5,
      "unit": "rpm",
      "quality": "GOOD"
    },
    {
      "name": "axis_x.motor_temp",
      "value": 72.1,
      "unit": "degC",
      "quality": "GOOD"
    }
  ]
}
\end{lstlisting}

Das Feld \texttt{quality} folgt der OPC-UA-Qualitätssemantik (GOOD, UNCERTAIN, BAD)
und erlaubt der App, fehlerhafte oder veraltete Werte visuell zu kennzeichnen.

\subsection{Variable Mapping Konfiguration (HVM)}

Die HMI Variable Map (HVM) ist eine XML-Konfigurationsdatei, die die Abbildung von
EtherCAT-PDO/SDO-Indizes auf benannte HMI-Variablen beschreibt. Sie definiert
Datentyp, Skalierung, Einheit, Zugriffsrechte und die erforderliche Mindestrolle für
jeden Parameter.

\begin{lstlisting}[language=xml, caption={HMI Variable Map (HVM) Konfigurationsbeispiel.}]
<HMIVariableMap xmlns="urn:ethercat:ehae:hvm:1.0">
  <Device slaveAddress="1001" name="Drive_Axis_X">
    <Variable
      hmiName="axis_x.actual_velocity"
      pdoIndex="0x6044" pdoSubIndex="0x00"
      dataType="INT16" scale="0.1" unit="rpm"
      access="READ" minRole="OPERATOR" />
    <Variable
      hmiName="axis_x.target_velocity"
      pdoIndex="0x60FF" pdoSubIndex="0x00"
      dataType="INT32" unit="rpm"
      access="READWRITE" minRole="OPERATOR" />
    <Variable
      hmiName="axis_x.kp_gain"
      sdoIndex="0x60F6" sdoSubIndex="0x01"
      dataType="UINT16"
      access="READWRITE" minRole="ENGINEER" />
    <Variable
      hmiName="axis_x.motor_temp"
      pdoIndex="0x2010" pdoSubIndex="0x03"
      dataType="INT16" scale="0.1" unit="C"
      access="READ" minRole="OPERATOR"
      alarmHigh="80.0" warnHigh="75.0" />
  </Device>
</HMIVariableMap>
\end{lstlisting}

Die Erweiterung um \texttt{alarmHigh}/\texttt{warnHigh}-Attribute erlaubt es dem Alarm
Manager, Grenzwertüberschreitungen direkt aus der HVM-Konfiguration abzuleiten, ohne
dass diese Logik in der App selbst implementiert werden muss.

% =====================================================================
\section{Sicherheitsarchitektur}
\label{sec:security}
% =====================================================================

Die Sicherheitsarchitektur der EHAE basiert auf drei aufeinander aufbauenden Säulen:
Transport-Verschlüsselung (TLS~1.3), Authentifizierung (OAuth~2.0 mit PKCE) und
Autorisierung (Role-Based Access Control, RBAC).

\subsection{Authentifizierungsablauf}

Abbildung~\ref{fig:auth_flow} zeigt die vollständige OAuth-2.0-PKCE-Sequenz für den
Verbindungsaufbau einer HMI-App zum HAG.

\begin{figure}[H]
\centering
\begin{tikzpicture}[
  actor/.style={rectangle, rounded corners=3pt, draw=etherblue, fill=etherblue!10,
    minimum width=2.5cm, minimum height=0.7cm, font=\footnotesize\bfseries,
    text=etherblue},
  msg/.style={-Stealth, thick},
  note/.style={font=\scriptsize, color=darkgray},
]
\node[actor] (app) at (0,0)  {HMI App};
\node[actor] (idp) at (5,0)  {Identity Provider};
\node[actor] (hag) at (10,0) {HAG};
\draw[dashed, color=darkgray] (0,-0.35)  -- (0,-9);
\draw[dashed, color=darkgray] (5,-0.35)  -- (5,-9);
\draw[dashed, color=darkgray] (10,-0.35) -- (10,-9);
\draw[msg] (0,-1)    -- (5,-1)
  node[midway, above, note] {1. Auth Request + code\_challenge (PKCE)};
\draw[msg] (5,-2)    -- (0,-2)
  node[midway, above, note] {2. Login-Seite};
\draw[msg] (0,-3)    -- (5,-3)
  node[midway, above, note] {3. Credentials + code\_verifier};
\draw[msg] (5,-4)    -- (0,-4)
  node[midway, above, note] {4. Access Token (JWT, 1h TTL) + Refresh Token};
\draw[msg] (0,-5)    -- (10,-5)
  node[midway, above, note] {5. WebSocket connect + Bearer Token (TLS 1.3)};
\draw[msg] (10,-6)   -- (5,-6)
  node[midway, above, note] {6. Token Introspection Request};
\draw[msg] (5,-7)    -- (10,-7)
  node[midway, above, note] {7. Token g\"{u}ltig, Rolle = OPERATOR};
\draw[msg] (10,-7.8) -- (0,-7.8)
  node[midway, above, note] {8. Verbindung akzeptiert, initiale Variable DB};
\draw[msg] (0,-8.8)  -- (10,-8.8)
  node[midway, above, note] {9. SUBSCRIBE \{axis\_x.*, conveyor.*\}};
\end{tikzpicture}
\caption{EHAE-Authentifizierungssequenz: OAuth~2.0 mit PKCE (Proof Key for Code
  Exchange, RFC~7636). Das Sequenzdiagramm zeigt drei Akteure auf vertikalen
  gestrichelten Lebenslinien: die \emph{HMI App} (links), den \emph{Identity Provider}
  (Mitte) und den \emph{HAG} (rechts). In Schritt~1 sendet die App einen Auth-Request
  mit einem PKCE-Code-Challenge an den Identity Provider. Dieser antwortet mit einer
  Login-Seite (Schritt~2). Nach Eingabe der Credentials und Übermittlung des Code-Verifiers
  (Schritt~3) stellt der Identity Provider ein JWT-Access-Token mit 1~Stunde Laufzeit und
  ein Refresh-Token aus (Schritt~4). In Schritt~5 verbindet sich die App über TLS~1.3
  mit dem HAG und übergibt das Bearer-Token. Der HAG verifiziert das Token durch
  Token-Introspection beim Identity Provider (Schritte~6--7) und akzeptiert die Verbindung
  mit Übertragung der initialen Variable-Datenbank (Schritt~8). Abschlie\ss{}end
  abonniert die App ihre benötigten Variablen (Schritt~9).}
\label{fig:auth_flow}
\end{figure}

Die Verwendung von PKCE \cite{rfc7636} ist entscheidend für die Sicherheit auf mobilen
Geräten, da diese keine Client-Secrets sicher speichern können. PKCE stellt sicher, dass
nur der Initiator des Auth-Flows das Token einlösen kann, selbst wenn der Authorization
Code abgefangen wird.

\subsection{Token-Management}

JWTs enthalten folgende standardisierten Claims gemä\ss{} RFC~7519:

\begin{itemize}[leftmargin=2em]
  \item \textbf{sub:} Eindeutige Benutzer-ID.
  \item \textbf{iat/exp:} Ausstellungs- und Ablaufzeit (1 Stunde TTL).
  \item \textbf{ehae\_role:} EHAE-spezifischer Claim für die Rolle (VIEWER, OPERATOR,
    ENGINEER, ADMIN).
  \item \textbf{ehae\_machines:} Whitelist der Maschinen-IDs, auf die dieser Token
    Zugriff gewährt.
\end{itemize}

Nach Ablauf des Access Tokens verwendet die App das Refresh Token, um still im
Hintergrund ein neues Access Token zu erhalten, ohne den Bediener mit einem erneuten
Login zu unterbrechen.

\subsection{Rollenbasierte Zugriffskontrolle}

Tabelle~\ref{tab:rbac} zeigt die vollständige RBAC-Matrix für alle vier EHAE-Rollen.

\begin{table}[H]
\centering
\caption{EHAE-RBAC-Matrix.}
\label{tab:rbac}
\begin{tabularx}{\textwidth}{lXXXX}
\toprule
\textbf{Fähigkeit}                & \textbf{Viewer} & \textbf{Operator} & \textbf{Engineer} & \textbf{Admin} \\
\midrule
Prozessdaten lesen                & \checkmark & \checkmark & \checkmark & \checkmark \\
Alarme quittieren                 &            & \checkmark & \checkmark & \checkmark \\
Sollwerte schreiben (PDO)         &            & \checkmark & \checkmark & \checkmark \\
Rezepte laden                     &            & \checkmark & \checkmark & \checkmark \\
SDO-Parameter lesen/schreiben     &            &            & \checkmark & \checkmark \\
Firmware-Update (FoE)             &            &            & \checkmark & \checkmark \\
Oszilloskop-Stream                &            &            & \checkmark & \checkmark \\
Benutzerverwaltung                &            &            &            & \checkmark \\
Rollenänderung                    &            &            &            & \checkmark \\
\bottomrule
\end{tabularx}
\end{table}

Die RBAC-Engine prüft bei jeder eingehenden \texttt{WRITE\_VAR}- oder
\texttt{SDO\_WRITE}-Nachricht sowohl die Rolle des Clients als auch die
\texttt{minRole}-Definition der jeweiligen Variable in der HVM-Konfiguration. Eine
Schreiboperation wird nur ausgeführt, wenn beide Bedingungen erfüllt sind.

\subsection{Netzwerksicherheit}

Alle WebSocket-Verbindungen sind obligatorisch über TLS~1.3 \cite{rfc8446} gesichert.
TLS~1.3 bietet gegenüber TLS~1.2 wichtige Verbesserungen:

\begin{itemize}[leftmargin=2em]
  \item Reduzierter Handshake auf 1 Round-Trip-Time (1-RTT).
  \item Entfernung unsicherer Cipher Suites (RC4, 3DES, etc.).
  \item Forward Secrecy für alle Verbindungen durch ephemere Diffie-Hellman-Schlüssel.
  \item 0-RTT-Wiederverbindung für bekannte Server (mit Replay-Schutz).
\end{itemize}

% =====================================================================
\section{Netzwerkintegration}
\label{sec:network}
% =====================================================================

\begin{figure}[h]
\centering
\begin{tikzpicture}[
  device/.style={rectangle, rounded corners=4pt, draw=#1, fill=#1!12,
    minimum width=2.2cm, minimum height=0.8cm, font=\footnotesize\bfseries, text=#1,
    align=center},
  sw/.style={rectangle, rounded corners=2pt, draw=darkgray, fill=lightgray,
    minimum width=2cm, minimum height=0.6cm, font=\scriptsize},
  cloud/.style={ellipse, draw=cloudblue, fill=cloudblue!15, minimum width=2.5cm,
    minimum height=1cm, font=\footnotesize, text=cloudblue},
  arr/.style={-Stealth, thick},
]
\node[cloud]          (cloud)   at (7,5)    {Cloud Monitor};
\node[device=darkgray](fw)      at (7,3)    {Firewall\\Router};
\node[sw]             (sw_corp) at (3.5,1.5){Corp VLAN 10};
\node[sw]             (sw_prod) at (7,1.5)  {Prod VLAN 20};
\node[sw]             (sw_hmi)  at (10.5,1.5){HMI VLAN 30};
\node[device=etherblue](master) at (7,-0.3) {EtherCAT\\Master + HAG};
\node[device=tabletblue](ap)    at (10.5,-0.3){Wi-Fi AP\\802.11ax};
\node[device=tabletblue](t1)    at (9,-2)   {Tablet\\Operator};
\node[device=appgreen] (t2)     at (11.5,-2){Laptop\\Engineer};
\node[device=darkgray] (pc)     at (3.5,-0.3){Office PC\\Viewer};
\node[device=etherblue](slaves) at (7,-2)   {EtherCAT\\Slaves};
\draw[arr] (cloud)   -- (fw);
\draw[arr] (fw)      -- (sw_corp);
\draw[arr] (fw)      -- (sw_prod);
\draw[arr] (fw)      -- (sw_hmi);
\draw[arr] (sw_prod) -- (master);
\draw[arr] (sw_hmi)  -- (ap);
\draw[arr] (sw_corp) -- (pc);
\draw[arr, dashed, color=tabletblue] (ap) -- (t1);
\draw[arr, dashed, color=appgreen]   (ap) -- (t2);
\draw[arr, color=etherblue] (master) -- (slaves)
  node[midway, left, font=\tiny] {EtherCAT};
\end{tikzpicture}
\caption{Empfohlene Netzwerksegmentierung für EHAE mit drei VLANs. Das Diagramm
  zeigt die vollständige Netzwerktopologie von oben nach unten. Ganz oben steht der
  \emph{Cloud Monitor} (blau, Ellipse), der über einen zentralen \emph{Firewall/Router}
  (grau) mit drei separaten VLANs verbunden ist: VLAN~10 für das
  Unternehmens-/Büronetz (links), VLAN~20 für das Produktionsnetz (Mitte) und
  VLAN~30 für das HMI-Netz (rechts). Im Produktionsnetz (VLAN~20) befindet sich der
  \emph{EtherCAT Master mit HAG} sowie die direkt angebundenen
  \emph{EtherCAT Slaves}. Im HMI-Netz (VLAN~30) steht ein \emph{Wi-Fi Access Point
  802.11ax}, über den \emph{Operator-Tablet} und \emph{Engineer-Laptop}
  (gestrichelte Pfeile = WLAN) verbunden sind. Im Büronetz (VLAN~10) ist ein
  \emph{Office PC} für die Viewer-Rolle angebunden. Die Firewall kontrolliert den
  Zugriff zwischen den VLANs und zur Cloud.}
\label{fig:network_diagram}
\end{figure}

Die VLAN-Segmentierung ist eine zentrale Sicherheitsma\ss{}nahme: Sie verhindert, dass
kompromittierte HMI-Geräte direkten Zugriff auf das EtherCAT-Produktionsnetz erhalten.
Der gesamte Datenfluss läuft ausschlie\ss{}lich über den HAG, der als einziger
Kommunikationspunkt zwischen App-Clients und EtherCAT-Netz fungiert.

\subsection{Firewall-Regelwerk}

Das empfohlene Firewall-Regelwerk erlaubt folgende Verbindungen:

\begin{itemize}[leftmargin=2em]
  \item \textbf{VLAN~30 $\to$ VLAN~20, Port 8443:} WebSocket-Verbindungen der HMI-Apps
    zum HAG (TCP, TLS~1.3).
  \item \textbf{VLAN~10 $\to$ VLAN~20, Port 8443:} Viewer-Zugriff vom Office-PC.
  \item \textbf{VLAN~20 $\to$ Internet, Port 443:} Token-Introspection zum Cloud-IdP.
  \item \textbf{Alle anderen:} Default DENY.
\end{itemize}

\subsection{Latenzbudget}

Tabelle~\ref{tab:latency} zeigt das End-to-End-Latenzbudget für die app-basierte HMI.

\begin{table}[H]
\centering
\caption{End-to-End-Latenzbudget für app-basiertes HMI.}
\label{tab:latency}
\begin{tabularx}{\textwidth}{lXrr}
\toprule
\textbf{Stufe}         & \textbf{Beschreibung}          & \textbf{Typisch}             & \textbf{Maximum} \\
\midrule
EtherCAT-Zyklus        & PDO-Datenaktualisierung        & $250\,\si{\micro\second}$    & 1\,ms \\
HAG-Verarbeitung       & Mapping, Änderungserkennung    & 0{,}1\,ms                   & 0{,}5\,ms \\
WebSocket-Framing      & JSON-Serialisierung            & 0{,}5\,ms                   & 2\,ms \\
Wi-Fi~6                & Drahtlose Übertragung          & 1--3\,ms                    & 10\,ms \\
App-Rendering          & UI-Aktualisierung              & 8--16\,ms                   & 33\,ms \\
\midrule
\textbf{Gesamt E2E}    & Prozessdaten bis Display       & $\approx$10\,ms             & $\approx$47\,ms \\
\bottomrule
\end{tabularx}
\end{table}

Eine Gesamtlatenz von typisch 10\,ms ist für alle HMI-Anwendungsszenarien vollkommen
ausreichend. Die menschliche Wahrnehmungsschwelle für visuelle Reaktionen liegt bei ca.
100\,ms; damit bietet EHAE einen komfortablen 10-fachen Sicherheitsabstand. Lediglich
für hochdynamische Steuerungseingriffe (Not-Halt etc.) muss beachtet werden, dass
Not-Halt-Funktionen weiterhin direkt am Feldbus oder als Safety-Funktion (STO) im Antrieb
implementiert werden müssen --- nicht über die HMI-App.

%% =====================================================================
%\section{Operator-App UI-Referenz}
%\label{sec:ui}
%% =====================================================================
%
%\begin{figure}[H]
%\centering
%\begin{tikzpicture}[
%  screen/.style={rounded corners=10pt, draw=tabletblue, fill=white,
%    minimum width=7cm, minimum height=12cm},
%  bar/.style={rectangle, draw=none, fill=#1, minimum width=6.6cm, minimum height=0.7cm},
%  btn/.style={rounded corners=5pt, draw=#1, fill=#1!20,
%    minimum width=2.8cm, minimum height=0.9cm, font=\scriptsize\bfseries, text=#1},
%  card/.style={rounded corners=4pt, draw=darkgray!40, fill=lightgray,
%    minimum width=6.6cm, minimum height=1.2cm},
%]
%\node[screen] (scr) at (0,0) {};
%\node[bar=etherblue] at (0, 5.25) {};
%\node[font=\tiny\bfseries, text=white] at (0, 5.25) {Produktionslinie A | Schicht 2 | 14:32};
%\node[card] at (0, 4.2) {};
%\node[font=\small\bfseries, color=appgreen]  at (-1.5, 4.4) {LÄUFT};
%\node[font=\tiny]                            at (0, 4.0)    {Drehzahl: 1250 rpm | Modus: AUTO};
%\node[card] at (0, 2.8) {};
%\node[font=\scriptsize, color=warnorange]    at (-1.5, 2.95) {\textbf{WARNUNG} Temp Achse Z};
%\node[font=\tiny]                            at (-1.5, 2.65) {72\,\textdegree C / Grenzwert: 80\,\textdegree C};
%\node[btn=warnorange]                        at (2.3, 2.8)   {ACK};
%\node[card] at (0, 1.5) {};
%\node[font=\scriptsize\bfseries]             at (-1.8, 1.65) {Teile heute};
%\node[font=\large\bfseries, color=etherblue] at (-0.5, 1.45) {1\,284};
%\node[font=\scriptsize\bfseries]             at (1.8, 1.65)  {Ziel};
%\node[font=\large\bfseries, color=darkgray]  at (2.5, 1.45)  {1\,500};
%\node[card] at (0, 0.2) {};
%\node[font=\scriptsize]                      at (0, 0.45)   {Bandgeschwindigkeit};
%\fill[lightgray!50] (-2.8, 0.05) rectangle (2.8, 0.25);
%\fill[appgreen!60]  (-2.8, 0.05) rectangle (1.2, 0.25);
%\node[font=\tiny]                            at (0, -0.1)   {0{,}45 m/s (max: 1{,}2 m/s)};
%\node[btn=accentred] at (-1.6,-1.2) {STOP};
%\node[btn=appgreen]  at ( 1.6,-1.2) {START};
%\node[bar=lightgray] at (0, -5.25) {};
%\node[font=\tiny] at (-2.2,-5.25) {Home};
%\node[font=\tiny] at (-0.7,-5.25) {Trends};
%\node[font=\tiny] at ( 0.7,-5.25) {Alarme};
%\node[font=\tiny] at ( 2.2,-5.25) {Einstellungen};
%\draw[tabletblue, thick, rounded corners=10pt] (-3.5,-6) rectangle (3.5,6);
%\end{tikzpicture}
%\caption{Referenz-Layout: Operator-App Hauptansicht (Smartphone-Darstellung). Das
%  Mockup zeigt ein hochformatiges Smartphone-Display mit blauem Header-Bereich, der
%  Produktionslinie, Schichtbezeichnung und Uhrzeit anzeigt. Darunter folgen vier
%  Informationskarten: (1)~Maschinenstatus in Grün ("{}LÄUFT"{}) mit aktueller
%  Drehzahl 1250~rpm im Automatikmodus, (2)~eine orange Warnmeldung für Achse~Z mit
%  Temperatur 72\,\textdegree C bei einem Grenzwert von 80\,\textdegree C sowie einem "{}ACK"{}-Button zum
%  Quittieren, (3)~ein Produktionszähler mit aktuellem Stand 1\,284 Teilen gegenüber
%  dem Tagesziel von 1\,500, (4)~ein Fortschrittsbalken für die Bandgeschwindigkeit
%  (0{,}45~m/s von max. 1{,}2~m/s, etwa 37\,\% der maximalen Geschwindigkeit).
%  Am unteren Rand befinden sich prominente "{}STOP"{}- und "{}START"{}-Buttons
%  sowie eine Navigation mit vier Reitern (Home, Trends, Alarme, Einstellungen).}
%\label{fig:operator_ui}
%\end{figure}
%
%Das UI-Design folgt den Prinzipien des \emph{Progressive Disclosure}: die
%wichtigsten Informationen (Status, kritische Alarme) sind sofort sichtbar, während
%detailliertere Ansichten (Trends, Parameterdetails) durch Navigation in weiteren
%Reitern erreichbar sind. Die Farbcodierung ist konsistent: Grün für Normalbetrieb,
%Orange für Warnungen, Rot für Fehler und Notoperationen.

% =====================================================================
\section{Migrationsstrategie}
\label{sec:migration}
% =====================================================================

\begin{figure}[h]
\centering
\begin{tikzpicture}[
  phase/.style={rectangle, rounded corners=5pt, draw=#1, fill=#1!15,
    minimum width=2.8cm, minimum height=2.5cm, font=\scriptsize, align=center, text=#1},
  arr/.style={-Stealth, very thick, color=darkgray},
]
\node[phase=etherblue] (p1) at (0,0) {\textbf{Phase 1}\\[4pt]Assessment\\[3pt]
  $\bullet$ Panels inventarisieren\\$\bullet$ PDO/SDO-Anforderungen\\$\bullet$ Rollen definieren};
\node[phase=appgreen]  (p2) at (4,0) {\textbf{Phase 2}\\[4pt]Parallelbetrieb\\[3pt]
  $\bullet$ HAG deployen\\$\bullet$ Mapping validieren\\$\bullet$ Schulung};
\node[phase=warnorange](p3) at (8,0) {\textbf{Phase 3}\\[4pt]App-\\Rollout\\[3pt]
  $\bullet$ Tablets verteilen\\$\bullet$ Engineer App\\$\bullet$ Panels entfernen};
\node[phase=darkgray]  (p4) at (12,0){\textbf{Phase 4}\\[4pt]Vollbetrieb\\EHAE\\[3pt]
  $\bullet$ 100\% app-basiert\\$\bullet$ Remote-Monitoring\\$\bullet$ Updates};
\draw[arr] (p1.east) -- (p2.west);
\draw[arr] (p2.east) -- (p3.west);
\draw[arr] (p3.east) -- (p4.west);
\draw[thick, color=darkgray] (-1,-1.8) -- (13,-1.8);
\foreach \x/\t in {0/0--1 Mo., 4/1--3 Mo., 8/3--6 Mo., 12/6+ Mo.} {
  \fill[darkgray] (\x,-1.8) circle (3pt);
  \node[font=\tiny, color=darkgray] at (\x,-2.2) {\t};
}
\end{tikzpicture}
\caption{Vierstufige Migrationsstrategie von traditionellen HMI-Panels zu EHAE. Das
  Phasendiagramm zeigt vier sequentielle Phasen auf einer Zeitachse. \textbf{Phase~1
  (blau, 0--1 Monat):} Assessment --- alle vorhandenen HMI-Panels werden inventarisiert,
  PDO/SDO-Anforderungen analysiert und Benutzerrollen definiert. \textbf{Phase~2 (grün,
  1--3 Monate):} Parallelbetrieb --- der HAG wird deployt und parallel zum bestehenden
  HMI-Panel betrieben; das Variable-Mapping wird validiert und Mitarbeiter werden
  geschult. \textbf{Phase~3 (orange, 3--6 Monate):} App-Rollout --- Tablets werden an
  Operatoren verteilt, die Engineer App wird für alle Inbetriebnahmen genutzt, und die
  traditionellen Panels werden schrittweise entfernt. \textbf{Phase~4 (grau, ab 6
  Monate):} Vollbetrieb EHAE --- 100\,\% app-basierter Betrieb mit Remote-Monitoring
  und Over-the-Air-Updates.}
\label{fig:migration_phases}
\end{figure}

Die Migrationsstrategie ist so konzipiert, dass zu keinem Zeitpunkt ein Produktionsausfall
entsteht. In Phase~2 laufen beide Systeme parallel, und erst nach vollständiger Validierung
des app-basierten HMI werden die Panels abgebaut. Dies minimiert das Risiko und erlaubt
einen kontrollierten Übergang auch in laufenden Produktionsumgebungen.

\subsubsection{Risikominderung}

Folgende Ma\ss{}nahmen reduzieren das Migrationsrisiko:

\begin{itemize}[leftmargin=2em]
  \item \textbf{Fallback-Plan:} Während Phase~2 und~3 bleibt das traditionelle Panel als
    Rückfalloption aktiv.
  \item \textbf{Pilotmaschine:} Die Migration beginnt an einer unkritischen Pilotmaschine,
    bevor sie auf die gesamte Produktionslinie ausgerollt wird.
  \item \textbf{Schrittweise Rollenfreigabe:} Operators erhalten zunächst nur
    Lesezugriff; Schreibrechte werden nach Schulungsabschluss freigeschaltet.
  \item \textbf{Audit-Log:} Alle Schreiboperationen und Alarmereignisse werden im HAG
    protokolliert und sind für Audits nachverfolgbar.
\end{itemize}

% =====================================================================
\section{Mobile App: Android-Implementierung mit Capacitor}
\label{sec:mobile}
% =====================================================================

Der EHAE-Referenz-Implementierung liegt eine entscheidende Architekturentscheidung
zugrunde: Die Mobile App für Android ist \emph{keine} eigenständige Entwicklung,
sondern eine direkte Kapselung der bereits vollständig funktionsfähigen Web-App in eine
native Android-Anwendung --- realisiert durch das \textbf{Capacitor-Framework}
\cite{capacitor} von Ionic. Dieses Prinzip der \emph{web-getriebenen
Mobile-App-Entwicklung} maximiert den Code-Reuse und minimiert den Wartungsaufwand
erheblich: jede Änderung an der Web-App steht automatisch auch in der Mobile App zur
Verfügung, ohne dass eine zweite Codebase gepflegt werden muss.

\subsection{Das Capacitor-Prinzip}

Capacitor ist ein plattformübergreifendes App-Laufzeit-Framework, das Web-Apps
(HTML, CSS, JavaScript) als vollwertige native Apps auf Android und iOS ausführt.
Im Gegensatz zu einem klassischen Hybrid-App-Ansatz (z.\,B.\ Cordova/PhoneGap)
bettet Capacitor die Web-App in eine native \texttt{WebView}-Komponente
(\texttt{android.webkit.WebView}) ein und stellt eine JavaScript-Bridge zur Verfügung,
über die Web-Code auf native Geräteschnittstellen zugreifen kann.

\begin{figure}[h]
\centering
\begin{tikzpicture}[
  layer/.style={rectangle, rounded corners=5pt, draw=#1, fill=#1!12,
    minimum width=11cm, minimum height=0.9cm, font=\small\bfseries, text=#1},
  arr/.style={-Stealth, thick},
]
\node[layer=darkgray]   (android) at (0,0)   {Android OS (AOSP)};
\node[layer=etherblue]  (shell)   at (0,1.4) {Capacitor Native Shell --- \texttt{CapacitorActivity} (Java/Kotlin)};
\node[layer=warnorange] (bridge)  at (0,2.8) {Capacitor JavaScript Bridge --- native $\leftrightarrow$ web};
\node[layer=appgreen]   (wv)      at (0,4.2) {Android WebView (\texttt{Chromium})};
\node[layer=tabletblue] (webapp)  at (0,5.6) {\textbf{EHAE Web-App} (HTML5 / CSS3 / JavaScript)};
\draw[arr, color=etherblue,  <->] (android.north) -- (shell.south);
\draw[arr, color=warnorange, <->] (shell.north)   -- (bridge.south);
\draw[arr, color=appgreen,   <->] (bridge.north)  -- (wv.south);
\draw[arr, color=tabletblue, <->] (wv.north)      -- (webapp.south);
\draw[decorate, decoration={brace,amplitude=8pt}, color=accentred]
  (6.4,0.5) -- (6.4,2.3) node[midway, right=8pt, font=\tiny, color=accentred, align=left]
  {Native\\Android};
\draw[decorate, decoration={brace,amplitude=8pt}, color=appgreen]
  (6.4,3.3) -- (6.4,6.1) node[midway, right=8pt, font=\tiny, color=appgreen, align=left]
  {Web-Layer\\(plattform-\\unabhängig)};
\end{tikzpicture}
\caption{Capacitor-Schichtenmodell der EHAE-Android-App. Das Diagramm zeigt fünf
  aufeinander aufbauende Schichten von unten nach oben. Ganz unten (grau) das
  \emph{Android-Betriebssystem (AOSP)}. Darüber (blau) die \emph{Capacitor Native
  Shell}: eine \texttt{CapacitorActivity}, die den Android-Lebenszyklus verwaltet.
  Die dritte Schicht (orange) ist die \emph{Capacitor JavaScript Bridge}: sie
  ermöglicht den bidirektionalen Datenaustausch zwischen nativem Kotlin/Java-Code
  und dem JavaScript der Web-App. Die vierte Schicht (grün) ist die
  \emph{Android WebView} auf Chromium-Basis, die die Web-App rendert.
  Die oberste Schicht (blau) ist die eigentliche \emph{EHAE Web-App} --- dieselbe
  \texttt{index.html}, die auch als Browser-Anwendung genutzt wird. Rechts zeigen
  zwei geschweifte Klammern die Trennung zwischen nativem Android-Layer und dem
  plattformunabhängigen Web-Layer.}
\label{fig:capacitor_stack}
\end{figure}

Die entscheidende Konsequenz dieses Modells: Die gesamte EHAE-Applikationslogik
--- Dashboard, Topologie-Viewer, Alarm-Manager, Oszilloskop, SDO-Browser,
TCO-Kostenanalyse --- ist \emph{ausschlie\ss{}lich} in der Web-App implementiert
und muss nur \emph{einmal} entwickelt und gewartet werden. Capacitor fungiert dabei
als dünner Wrapper, der die Web-App in eine vollwertige Android-APK verwandelt.

\subsection{Projektstruktur}

Die Android-App-Implementierung liegt unter \texttt{src/android/ehae-android/} und
folgt der Standard-Capacitor-Projektstruktur:

\begin{lstlisting}[language={}, caption={Projektstruktur der EHAE-Android-App.}]
ehae-android/
|-- www/                    # Web-Assets (EHAE HMI SPA)
|   `-- index.html          # Vollstaendige EHAE Web-App
|-- android/                # Natives Android-Gradle-Projekt
|   |-- app/
|   |   |-- src/main/       # Android Activity, XML-Manifest
|   |   `-- build.gradle    # App-Level Build-Konfiguration
|   |-- capacitor-cordova-android-plugins/
|   `-- build.gradle        # Projekt-Level Build-Konfiguration
|-- capacitor.config.json   # Capacitor-Hauptkonfiguration
`-- package.json            # npm-Abhaengigkeiten
\end{lstlisting}

Das Verzeichnis \texttt{www/} enthält die Web-App-Assets --- identisch mit der
Desktop-Browser-Version unter \texttt{src/index.html}. Capacitor liest diese Assets
beim Build-Vorgang ein und bündelt sie in das Android-App-Bundle, so dass die App
offline und ohne externe Netzwerkzugriffe auf die eigene UI-Logik lauffähig ist.

\subsection{Capacitor-Konfiguration}

Die zentrale Konfigurationsdatei \texttt{capacitor.config.json} steuert das Verhalten
der nativen Wrapper-Schicht:

\begin{lstlisting}[language={}, caption={Capacitor-Konfiguration (\texttt{capacitor.config.json}).}]
{
  "appId": "de.epp.ehae",
  "appName": "EHAE HMI",
  "webDir": "www",
  "server": { "androidScheme": "https" },
  "android": {
    "allowMixedContent": true,
    "backgroundColor": "#07090f",
    "overrideUserAgent": "EHAE-HMI/1.0 Android"
  },
  "plugins": {
    "SplashScreen": {
      "launchShowDuration": 1500,
      "backgroundColor": "#07090f",
      "showSpinner": false
    }
  }
}
\end{lstlisting}

Die wichtigsten Konfigurationsparameter im Einzelnen:

\begin{description}[leftmargin=2em]
  \item[\texttt{appId}] Eindeutiger Android-App-Bezeichner nach Reverse-Domain-Notation
    (\texttt{de.epp.ehae}), der die App im Google Play Store und auf dem Gerät identifiziert.
  \item[\texttt{webDir}] Verzeichnis mit den Web-Assets. Der Inhalt des
    \texttt{www/}-Verzeichnisses dient als Einstiegspunkt der WebView-Anwendung.
  \item[\texttt{androidScheme: "https"}] Erzwingt das HTTPS-Schema für interne
    WebView-URLs, damit der Browser-Sicherheitskontext und Cookies korrekt
    funktionieren.
  \item[\texttt{allowMixedContent}] Erlaubt gemischte HTTP/HTTPS-Verbindungen im
    Intranet-Kontext, damit die App auch über unverschlüsselte lokale HAG-Verbindungen
    kommunizieren kann.
  \item[\texttt{backgroundColor}] Setzt die initiale Hintergrundfarbe auf das
    EHAE-Dunkelblau (\texttt{\#07090f}) für einen nahtlosen Übergang vom Splash-Screen
    zur vollständig geladenen Web-App.
  \item[\texttt{overrideUserAgent}] Identifiziert App-Anfragen als
    \texttt{EHAE-HMI/1.0 Android}, was im HAG für Logging und
    gerätetyp-spezifische Optimierungen genutzt werden kann.
\end{description}

\subsection{Abhängigkeiten und Framework-Version}

Tabelle~\ref{tab:npm_deps} listet die npm-Abhängigkeiten der Android-App auf.

\begin{table}[h]
\centering
\caption{npm-Abhängigkeiten der EHAE-Android-App (\texttt{package.json}).}
\label{tab:npm_deps}
\begin{tabularx}{\textwidth}{llX}
\toprule
\textbf{Paket} & \textbf{Version} & \textbf{Funktion} \\
\midrule
\texttt{@capacitor/core}    & \textasciicircum{}6.2.1 & Capacitor-Kern-Runtime, JavaScript-Bridge \\
\texttt{@capacitor/android} & \textasciicircum{}6.2.1 & Native Android-Implementie-\ rung \\
\texttt{@capacitor/cli}     & \textasciicircum{}6.2.1 & Build- und Sync-Werkzeuge (\texttt{cap sync}) \\
\parbox[t]{5cm}{\raggedright\texttt{@capawesome/capacitor-}\\\texttt{android-edge-to-edge-support}}
                            & \textasciicircum{}8.0.6 & Vollbild-Edge-to-Edge-Darstellung \\
\bottomrule
\end{tabularx}
\end{table}

Die Verwendung von Capacitor~6 stellt sicher, dass die Web-App auf Android~5.1 bis
Android~14 (\texttt{minSdkVersion}~22, \texttt{targetSdkVersion}~34) lauffähig ist.
Das Plugin \texttt{@capawesome/\newline capacitor-android-edge-to-edge-support} ermöglicht die
vollkantige Darstellung ohne Statusbar- oder Navigationsbar-Unterbrechung --- kritisch
für HMI-Dashboards, die maximale Bildschirmfläche benötigen.

\subsection{Build-Workflow}

Abbildung~\ref{fig:capacitor_build} zeigt den Build-Workflow von der Web-App-Entwicklung
bis zur verteilbaren Android-APK.

\begin{figure}[h]
\centering
\begin{tikzpicture}[
  step/.style={rectangle, rounded corners=4pt, draw=#1, fill=#1!15,
    minimum width=2.8cm, minimum height=1.0cm, font=\footnotesize\bfseries,
    text=#1, align=center},
  arr/.style={-Stealth, very thick, color=darkgray},
  note/.style={font=\tiny, color=darkgray, align=center},
]
\node[step=tabletblue] (web)   at (0,0)    {Web-App\\entwickeln};
\node[step=etherblue]  (copy)  at (3.5,0)  {Assets in\\www/ kopieren};
\node[step=appgreen]   (sync)  at (7,0)    {\texttt{cap sync}\\android};
\node[step=warnorange] (build) at (10.5,0) {\texttt{./gradlew}\\assembleRelease};
\node[step=accentred]  (apk)   at (10.5,-2.2) {APK / AAB};
\draw[arr] (web.east)   -- (copy.west);
\draw[arr] (copy.east)  -- (sync.west);
\draw[arr] (sync.east)  -- (build.west);
\draw[arr] (build.south) -- (apk.north);
\node[note] at (1.75,-0.7) {\texttt{src/index.html}};
\node[note] at (5.25,-0.7) {Web-Assets in\\Gradle-Projekt};
\node[note] at (9.0, -0.7) {Gradle-Build};
\end{tikzpicture}
\caption{Build-Workflow der EHAE-Android-App. Das Ablaufdiagramm zeigt vier Schritte.
  \textbf{Schritt~1 (blau):} Die Web-App wird entwickelt und direkt im Browser
  getestet --- ohne Android-Abhängigkeiten. \textbf{Schritt~2 (blau):} Die fertigen
  Web-Assets (\texttt{index.html}) werden in das \texttt{www/}-Verzeichnis des
  Capacitor-Projekts kopiert. \textbf{Schritt~3 (grün):} \texttt{cap sync android}
  synchronisiert die Web-Assets und alle Capacitor-Plugins in das native
  Android-Gradle-Projekt. \textbf{Schritt~4 (orange):} \texttt{./gradlew
  assembleRelease} baut das vollständige Android-App-Bundle. Das Ergebnis (rot) ist
  eine APK oder ein AAB, die direkt auf Android-Geräten installiert oder im
  Google Play Store veröffentlicht werden kann.}
\label{fig:capacitor_build}
\end{figure}

Der Build-Prozess in der Kommandozeile:

\begin{lstlisting}[language=bash, caption={EHAE-Android-App: Build-Ablauf.}]
# Schritt 1: Capacitor-Abhaengigkeiten installieren
npm install

# Schritt 2: Web-Assets synchronisieren
npx cap sync android

# Schritt 3: Debug-APK bauen und installieren
cd android && ./gradlew assembleDebug
adb install app/build/outputs/apk/debug/app-debug.apk

# Release-APK (fuer Produktion) bauen
./gradlew assembleRelease
\end{lstlisting}

\subsection{Vorteile der web-getriebenen Mobile-App-Entwicklung}

Tabelle~\ref{tab:capacitor_advantages} fasst die wichtigsten Vorteile des
Capacitor-Ansatzes gegenüber einer nativen Android-Entwicklung zusammen.

\begin{table}[H]
\centering
\caption{Vergleich: Capacitor (Web-Wrapper) vs.\ native Android-Entwicklung.}
\label{tab:capacitor_advantages}
\begin{tabularx}{\textwidth}{lXX}
\toprule
\textbf{Aspekt} & \textbf{Capacitor (Web-Wrapper)} & \textbf{Native Android-Entwicklung} \\
\midrule
Codebase           & Eine gemeinsame Codebase für Web und Android
                   & Separate Android-Codebase (Kotlin/Java) \\
Wartungsaufwand    & Einmalige Pflege der Web-App
                   & Doppelter Aufwand (Web und Android) \\
Feature-Parität    & Automatisch: Web $\equiv$ App
                   & Manuelle Synchronisation erforderlich \\
UI-Tests           & Direkt im Browser ausführbar
                   & Android-Emulator oder -Gerät notwendig \\
Update-Zyklus      & Web-App-Update $=$ App-Update
                   & Separater Gradle-Build und APK-Release \\
Zielplattformen    & Web, Android, iOS aus einer Codebase
                   & Ausschlie\ss{}lich Android \\
Einstiegshürde     & Niedrig (Web-Kenntnisse genügen)
                   & Hoch (Android-SDK-Kenntnisse erforderlich) \\
\bottomrule
\end{tabularx}
\end{table}

Der praktische Nutzen manifestiert sich im gesamten Entwicklungs-Workflow: Neue HMI-%
Features werden ausschlie\ss{}lich in der Web-App (\texttt{src/index.html}) entwickelt
und sofort im Browser getestet. Nach der Verifikation genügt ein \texttt{cap sync
android} gefolgt von einem Gradle-Build, um die identische Funktionalität als
Android-APK auszuliefern. Es entsteht kein redundanter Code, keine Divergenz zwischen
Web- und App-Version und kein zusätzlicher Qualitätssicherungsaufwand für eine zweite
Plattform. Dieses Prinzip spiegelt exakt das EHAE-Gesamtkonzept wider: maximale
Standardisierung und minimaler proprietärer Aufwand.

\subsection{Android-spezifische Anpassungen}

Trotz des Wrapper-Ansatzes sind einige Android-spezifische Anpassungen notwendig, um
eine native Nutzererfahrung zu bieten:

\begin{description}[leftmargin=2em]
  \item[Edge-to-Edge-Darstellung] Das Plugin
    \texttt{@capawesome/capacitor-android-edge-to\newline -edge-support} deaktiviert die
    Android-Systemleisten (Statusbar, Navigationsleiste) und erlaubt der EHAE-Web-App,
    das gesamte Displayformat für das HMI-Dashboard zu nutzen --- entscheidend für
    Querformat-HMI-Layouts auf Industrietablets.
  \item[Querformat-Fixierung] Das Android-Manifest konfiguriert
    \texttt{screenOrientation="{}land\newline scape"{}}, sodass die App immer im Querformat
    geöffnet wird. HMI-Dashboards sind auf Querformat optimiert und würden im
    Hochformat wichtige Informationsdichte verlieren.
  \item[Intranet-Verbindungen] Im Produktionsnetz kommuniziert die App über
    unverschlüsselte HTTP-Verbindungen mit dem lokalen HAG.
    \texttt{allowMixedContent: true} und \texttt{usesCleartextTraffic} im
    Android-Manifest erlauben diese Verbindungen im Intranet-Kontext, ohne die
    TLS-Anforderung für externe Verbindungen aufzuweichen.
  \item[Splash-Screen] Ein 1{,}5-sekündiger Splash-Screen in der EHAE-Hintergrundfarbe
    \texttt{\#07090f} überbrückt die WebView-Ladezeit und vermittelt ein
    professionelles, natives Erscheinungsbild.
  \item[User-Agent-Identifikation] Der überschriebene User-Agent
    (\texttt{EHAE-HMI/1.0 Android}) erlaubt dem HAG, Mobile-App-Verbindungen
    von Browser-Verbindungen zu unterscheiden.
\end{description}

% =====================================================================
\section{Evaluation}
\label{sec:eval}
% =====================================================================

\subsection{Kumulative Kostenanalyse}

Abbildung~\ref{fig:cumulative_cost} zeigt die kumulative 5-Jahres-Kostenentwicklung für
eine typische Anlage mit 20 HMI-Stationen.

\begin{figure}[h]
\centering
\begin{tikzpicture}
\begin{axis}[
  width=13cm, height=7cm,
  xlabel={Jahr}, ylabel={Kumulierte Kosten (EUR, Tausend)},
  xmin=0, xmax=5, ymin=0, ymax=220,
  xtick={0,1,2,3,4,5},
  legend pos=north west,
  ymajorgrids=true, grid style=dashed,
  title={Kumulierte 5-Jahres-Kosten: 20 HMI-Stationen},
  title style={font=\bfseries},
]
\addplot[color=etherblue, thick, mark=square*] coordinates {
  (0,0)(1,70)(2,110)(3,148)(4,178)(5,214)};
\addplot[color=appgreen, thick, mark=*] coordinates {
  (0,0)(1,16)(2,22)(3,28)(4,34)(5,40)};
\legend{Traditionelles HMI-Panel (20 Stationen), App-Based HMI (EHAE)}
\end{axis}
\end{tikzpicture}
\caption{Kumulative 5-Jahres-Kosten für 20 HMI-Stationen: traditionelles HMI-Panel vs.\
  EHAE-Lösung. Das Liniendiagramm zeigt zwei Kurven über fünf Jahre. Die \textbf{blaue
  Kurve} (traditionelle HMI-Panels) steigt steil an: nach Jahr~1 bereits \EUR{70\,000}
  (20 Panels $\times$ \EUR{3\,500}), nach Jahr~2 \EUR{110\,000} (Wartungskosten
  akkumulieren), und erreicht nach Jahr~5 \EUR{214\,000}. Die \textbf{grüne Kurve}
  (EHAE app-basiert) verläuft deutlich flacher: Jahr~1 \EUR{16\,000} (Tablets +
  HAG-Softwarelizenzen), danach nur geringe Zuwächse durch minimale Wartungskosten,
  Endstand nach Jahr~5: \EUR{40\,000}. Dies entspricht einer \textbf{Kostenreduktion
  von 81\,\%}. Die Schere zwischen beiden Kurven öffnet sich von Jahr zu Jahr weiter,
  da traditionelle Panels hohe laufende Wartungs- und Lizenzkosten verursachen.}
\label{fig:cumulative_cost}
\end{figure}

\subsection{Nicht-monetäre Vorteile}

Neben den direkten Kostenersparnissen bietet EHAE weitere Vorteile, die schwerer
monetär quantifizierbar, aber operativ signifikant sind:

\begin{description}[leftmargin=2em]
  \item[Verbesserte Diagnosefähigkeit] Die Engineer App mit integriertem Topologie-Viewer,
    Oszilloskop und SDO-Browser ermöglicht Inbetriebnahmen und Fehlerdiagnosen, die mit
    traditionellen Panels nicht möglich waren.
  \item[Ortsunabhängiges Monitoring] Fernzugriff über das Cloud-Monitoring-Modul erlaubt
    OEM-Support und Remote-Diagnose ohne Vor-Ort-Besuch.
  \item[Kürzere Update-Zyklen] App-Updates werden zentral über App-Stores ausgerollt;
    keine manuelle Panel-Aktualisierung notwendig.
  \item[Zukunftssicherheit] Consumer-Hardware-Ökosysteme (iOS, Android, Windows)
    werden langfristig von gro\ss{}en Unternehmen gepflegt und ständig verbessert.
\end{description}

% =====================================================================
\section{Schlussfolgerung}
\label{sec:conclusion}
% =====================================================================

Diese Arbeit hat die EtherCAT-HMI-App Extension (EHAE) vorgestellt, die
band-integrierte HMI-Controller aus EtherCAT-basierten Anlagen formal eliminiert. Die
zentralen Beiträge dieser Arbeit sind:

\begin{enumerate}[leftmargin=2em]
  \item Der \emph{HMI Application Gateway (HAG)}: ein neues Master-Stack-Modul, das
    keine Änderungen an der Slave-Firmware erfordert und vollständig rückwärtskompatibel
    ist.
  \item Das \emph{HAG Communication Protocol (HCP)}: WebSocket-basierter,
    JSON-serialisierter Echtzeit-Variablenzugriff mit Alarm-Delivery, SDO-Unterstützung,
    Oszilloskop-Strea\-ming und FoE-Firmware-Updates.
  \item Ein \emph{Dual-Role-App-Modell}: Operator App und Engineer App aus einer einzigen
    Codebase mit dynamischer Rollenaktivierung.
  \item Ein umfassendes Sicherheitsmodell auf Basis von OAuth~2.0 mit PKCE, TLS~1.3 und
    RBAC mit vier Rollen.
  \item Eine vierstufige Migrationsstrategie mit Fallback-Option und Risikominderungsplan.
  \item Nachgewiesene Kosteneinsparungen von bis zu 81\,\% über 5 Jahre bei gleichzeitiger
    erheblicher Verbesserung der Anzeigequalität und Diagnosefähigkeit.
\end{enumerate}

Wir laden die EtherCAT Technology Group ein, EHAE als formale Erweiterung des EtherCAT
Application Layer in die ETG.1000-Spezifikation aufzunehmen.

\subsection{Ausblick}

Die in dieser Arbeit vorgestellte EHAE-Architektur bildet eine tragfähige Basis für
eine Vielzahl weiterführender Entwicklungen. Die folgenden Abschnitte beschreiben
konkrete Forschungs- und Implementierungsrichtungen, die kurz-, mittel- und
langfristig verfolgt werden sollten, um das Potenzial app-basierter HMI-Clients im
industriellen EtherCAT-Umfeld vollständig auszuschöpfen.

\subsubsection*{Funktionale Sicherheit -- EHAE Safety Extension}

Das aktuelle EHAE-Modell ist für informationally-kritische, aber nicht für
sicherheitsgerichtete Anwendungen ausgelegt. Eine zukünftige \emph{EHAE Safety
Extension} würde einen separaten, redundant übertragenen HMI-Kanal nach
IEC~61508\,SIL\,2 und EN~ISO~13849\,PLd definieren. Dabei sind folgende
Teilprobleme zu lösen:

\begin{itemize}[leftmargin=2em]
  \item \textbf{Gesicherte Befehlsbestätigung:} Jeder sicherheitsrelevante
    Bedienerbefehl (z.\,B.\ Not-Halt-Quittierung, Schutzgitter-Freigabe) muss
    durch einen kryptografisch abgesicherten Handshake zwischen App und HAG
    quittiert werden. Das HCP-Protokoll ist um ein \verb|SAFE_ACK|-Feld zu
    ergänzen, das einen Hardware-Token (FIDO2/U2F) oder einen
    Biometrie-Nachweis des Bedieners einschließt.
  \item \textbf{Latenz-Budget-Garantien:} Für sicherheitsrelevante Kommandos
    muss die Ende-zu-Ende-Latenz (App $\to$ HAG $\to$ Slave) formal nachgewiesen
    werden. Reservierte Quality-of-Service-Klassen im WebSocket-Multiplexing und
    eine dedizierte TCP-Verbindung (alternativ QUIC-Datagram) sind vorzusehen.
  \item \textbf{Watchdog-Mechanismus:} Der HAG muss bei Verbindungsverlust zur
    Operator-App innerhalb einer konfigurierbaren Watchdog-Zeit ($t_{\text{WD}}
    \leq 500\,\si{\milli\second}$) autonom in einen sicheren Maschinenzustand
    übergehen und dies per \verb|SAFETY_TIMEOUT|-Alarm signalisieren.
  \item \textbf{Zertifizierungsweg:} Eine Zusammenarbeit mit einem Prüfinstitut
    (TÜV, DEKRA) ist notwendig, um eine typische EHAE-Safety-Installation nach
    IEC~62443-3-3 und EN~ISO~13849 zu zertifizieren. Das Ergebnis wäre ein
    standardisiertes Prüfverfahren für ETG-Mitglieder.
\end{itemize}

\subsubsection*{Drahtlose Echtzeit-Kommunikation -- TSN und Wi-Fi~7}

Die aktuelle EHAE-Spezifikation setzt auf kabelgebundenes Ethernet oder Standard-Wi-Fi.
Time-Sensitive Networking (TSN, IEEE~802.1Q-2022) und Wi-Fi~7 (IEEE~802.11be) eröffnen
neue Möglichkeiten:

\begin{itemize}[leftmargin=2em]
  \item \textbf{TSN-Integration:} Durch IEEE~802.1Qbv (Time-Aware Shaper) können
    HCP-Pakete in dedizierten Sendezeitfenstern (Guard Bands) priorisiert werden.
    Dies ermöglicht eine garantierte Latenz von unter $5\,\si{\milli\second}$ für
    PDO-Updates auch in überlasteten Netzwerken -- relevant für
    Montagelinien mit mehreren parallelen EHAE-Clients.
  \item \textbf{Wi-Fi~7 Multi-Link Operation (MLO):} IEEE~802.11be erlaubt
    simultane Übertragung über mehrere Frequenzbänder (2,4\,GHz, 5\,GHz, 6\,GHz).
    Für EHAE bedeutet dies eine drastische Reduktion des Jitters bei
    drahtlosen Tablets in der Fabrikhalle. Eine HAG-seitige Implementierung
    müsste differenzierte DSCP-Markierungen für HCP-Nachrichten gegenüber
    Hintergrundverkehr vergeben.
  \item \textbf{5G-Industrial-Private-Networks:} In weitläufigen Anlagen (z.\,B.\
    Logistikzentren) sind 5G-Campus-Netzwerke mit URLLC-Profil (\emph{Ultra
    Reliable Low Latency Communication}) eine realistische Option für
    EHAE-Clients auf fahrerlosen Transportsystemen (FTS). Der HAG muss in
    diesem Fall eine redundante Uplink-Umschaltung zwischen Wi-Fi und 5G
    unterstützen, ohne aktive WebSocket-Sessions zu unterbrechen (Session
    Migration via QUIC Connection ID Migration).
\end{itemize}

\subsubsection*{Künstliche Intelligenz und prädiktive Wartung}

Die PDO-Zeitreihendaten, die der HAG bereits sammelt, stellen ein hochwertiges
Rohmaterial für maschinelles Lernen dar. Konkrete Ausbaustufen:

\begin{itemize}[leftmargin=2em]
  \item \textbf{On-Device-Anomalieerkennung im HAG:} Ein leichtgewichtiges
    LSTM- oder Autoencoder-Modell (quantisiert auf INT8, z.\,B.\ via TensorFlow
    Lite) kann direkt auf dem HAG-Host-Prozessor laufen und statistische
    Ausreißer in Stromaufnahme, Temperatur oder Positionsfehlern erkennen,
    bevor ein Alarm ausgelöst wird. Die Schwellenwerte werden durch
    unüberwachtes Lernen in der Anlaufphase der Maschine automatisch kalibriert.
  \item \textbf{Cloud-basiertes Fleet-Learning:} Mehrere Anlagen desselben
    Typs können aggregierte (anonymisierte) Degradationsdaten in ein zentrales
    ML-Backend (z.\,B.\ Azure Machine Learning, AWS SageMaker) einspeisen.
    Das daraus gewonnene globale Modell wird per FoE über den HAG auf alle
    Slaves verteilt, die ein lokales Inferenz-Modul enthalten.
  \item \textbf{Natural Language Interface:} Mittelfristig ist ein
    KI-Assistent denkbar, der auf dem mobilen HMI-Client läuft und dem
    Bediener auf Basis aktueller Maschinendaten in natürlicher Sprache
    Diagnosen und Handlungsempfehlungen liefert (z.\,B.:\ \emph{``Achse 3
    zeigt erhöhten Schleppfehler -- Kugelgewindetrieb prüfen''}).
    Large Language Models (LLMs) mit On-Device-Inferenz auf
    Apple-Silicon- bzw.\ Qualcomm-Snapdragon-Chips wären hier praktikabel.
  \item \textbf{Reinforcement Learning für Rezeptoptimierung:} Für
    parametrierbare Prozesse (z.\,B.\ Schweißnähte, Dosierungen) könnte
    ein RL-Agent auf Basis beobachteter Qualitätskennzahlen die
    Maschinenparameter inkrementell optimieren -- mit dem HMI-Client als
    Supervision- und Override-Interface.
\end{itemize}

\subsubsection*{Augmented Reality und neue Interaktionsparadigmen}

Die EHAE-App-Schnittstelle ist für 2D-Touchscreens ausgelegt. Zukünftige
Interaktionsformen werden deutlich räumlicher:

\begin{itemize}[leftmargin=2em]
  \item \textbf{Spatial Computing (Apple Vision Pro, Meta Quest Pro):}
    Die WebXR Device API erlaubt Web-Apps, 3D-Overlays in das
    Kamerabild der Brille zu rendern -- ohne nativen Code. Da EHAE
    bereits auf einer Web-App basiert, genügt eine Erweiterung der
    HMI-Komponenten um \texttt{<a-frame>}-Elemente oder Three.js, um
    3D-Maschinensymbole über den physischen Slaves schweben zu lassen.
    PDO-Echtzeitwerte würden direkt am Slave-Gehäuse visualisiert.
  \item \textbf{Microsoft HoloLens 2 / RealWear HMT-1:} Für Wartungstechniker
    in industriellen Umgebungen ermöglicht eine EHAE-AR-Erweiterung
    hands-free Diagnose: QR-Code am Slave scannen, HAG liefert
    automatisch Topologie-Position, SDO-Werte und Fehlerliste im
    HoloLens-Overlay.
  \item \textbf{Gestensteuerung für sterile Umgebungen:} In der
    Pharma- und Lebensmittelproduktion, wo Touchscreens hygienisch
    problematisch sind, ermöglicht eine Integration von MediaPipe
    (Google) Hand Tracking den berührungslosen Betrieb des
    EHAE-HMI-Clients -- vollständig browserbasiert ohne
    zusätzliche Treiber.
  \item \textbf{Voice Control:} Die Web Speech API ist in allen
    modernen Browsern verfügbar. Sprachkommandos wie \emph{``Zeige
    Fehlerhistorie Achse 2''} oder \emph{``Starte Teach-Modus''} könnten
    maschinenlesbar in HCP-Nachrichten übersetzt werden,
    was den Bedienkomfort in Lärm- und Schutzhandschuh-Szenarien
    erheblich verbessert.
\end{itemize}

\subsubsection*{Edge Computing und HAG-Erweiterbarkeit}

Der HAG ist konzeptionell ein Middleware-Dienst. Seine Erweiterbarkeit zu einer
vollständigen Edge-Computing-Plattform ist ein natürlicher nächster Schritt:

\begin{itemize}[leftmargin=2em]
  \item \textbf{Plugin-Architektur:} Eine formale Plugin-API (z.\,B.\ auf Basis
    von WebAssembly-Modulen oder dynamisch geladenen Shared Libraries) würde
    es ETG-Mitgliedern erlauben, herstellerspezifische Protokoll-Adapter zu
    entwickeln (z.\,B.\ CoE-Profil-Decoder für SINAMICS-Antriebe, spezifische
    FoE-Handler). Das Plugin-System müsste hot-reload-fähig sein, um
    Produktionsstillstände bei Updates zu vermeiden.
  \item \textbf{Containerisierung:} Der HAG als Docker-Container (OCI-konform)
    auf einem industriellen Edge-Server (z.\,B.\ Siemens SIMATIC IPC, Dell
    Edge Gateway) ermöglicht orchestrierte Deployments via Kubernetes~K3s.
    Dies vereinfacht Rolling Updates, A/B-Deployments und Health-Monitoring.
  \item \textbf{Multi-Master-Unterstützung:} Anlagen mit mehreren EtherCAT-Linien
    (z.\,B.\ parallele Fertigungsstraßen) benötigen einen übergeordneten
    \emph{EHAE Aggregation Gateway}, der die HAG-Instanzen mehrerer Master
    unter einer einheitlichen HCP-Session zusammenführt. Der HMI-Client sieht
    dann eine einheitliche Topologie-Ansicht über alle Linien hinweg.
  \item \textbf{HAG-Redundanz (High Availability):} Für missionskritische Anlagen
    ist ein Active-Standby-HAG-Paar mit Shared State via Redis oder
    Apache Kafka denkbar. Failover unter $1\,\si{\second}$ ohne
    Client-Session-Unterbrechung wäre das Zielkriterium.
\end{itemize}

\subsubsection*{Digital Twin und IIoT-Integration}

Die strukturierten PDO- und SDO-Daten, die der HAG kontinuierlich verarbeitet,
sind ideale Eingaben für digitale Zwillinge und IIoT-Plattformen:

\begin{itemize}[leftmargin=2em]
  \item \textbf{OPC~UA~FX (Field eXchange):} Die neue OPC~UA~Part~80
    (\emph{Field eXchange}) spezifiziert ein Echtzeit-Kommunikationsprotokoll
    zwischen Feldgeräten über Hersteller- und Protokollgrenzen hinweg. Ein
    HAG-seitiger OPC~UA~FX-Publisher würde Ether\-CAT-PDOs als
    OPC~UA-Nodes exponieren und damit EHAE in industrielle IoT-Plattformen
    wie Azure Industrial IoT, Siemens MindSphere oder Bosch IoT Suite
    integrierbar machen.
  \item \textbf{Digital Twin nach ISO~23247:} Der ISO-Standard für digitale
    Zwillinge in der Fertigung definiert eine \emph{Observable Manufacturing
    Element (OME)}-Schnittstelle. Ein EHAE-Kompatibilitätsmodul könnte
    automatisch aus dem EtherCAT-ENI-File (EtherCAT Network Information) einen
    strukturierten Digital-Twin-Deskriptor generieren und diesen in
    DTDL (\emph{Digital Twin Definition Language}, Microsoft) oder
    AAS (\emph{Asset Administration Shell}, IEC~63278) ausdrücken.
  \item \textbf{MQTT-Sparkplug-B-Publisher:} Viele MES- und SCADA-Systeme
    erwarten MQTT-Sparkplug-B-Telemetrie. Ein HAG-Modul, das PDO-Werte
    automatisch als Sparkplug-B-Metriken unter einem konfigurierbaren
    Topic-Namespace publiziert, würde EHAE ohne weitere Middleware mit
    Systemen wie Ignition, Node-RED oder InfluxDB verbinden.
  \item \textbf{Historisierung und Langzeitarchivierung:} Ein integrierter
    Time-Series-Adapter (InfluxDB, TimescaleDB) im HAG würde PDO-Zeitreihen
    direkt archivieren -- mit konfigurierbarer Auflösung (z.\,B.\ $1\,\si{ms}$
    für Debugging, $1\,\si{s}$ für Langzeittrends) und automatischem Downsampling
    per Retention Policy. Die Engineer App könnte historische Daten im gleichen
    Oszilloskop-Interface darstellen wie Live-Daten.
\end{itemize}

\subsubsection*{App-Plattform-Erweiterung: iOS, Desktop und PWA}

Die aktuelle Implementierung fokussiert auf Android via Capacitor. Folgende
Plattformerweiterungen sind mit minimalem Aufwand realisierbar und betrieblich
relevant:

\begin{itemize}[leftmargin=2em]
  \item \textbf{iOS-App:} Capacitor unterstützt iOS nativ. Die vorhandene
    Web-App-Codebase kann ohne Anpassung für den Apple App Store kompiliert
    werden. Sicherheitsrelevante Ergänzungen für iOS wären die Integration von
    Face~ID/Touch~ID als PKCE-ergänzendem Second Factor sowie die Nutzung
    des iOS Secure Enclave für JWT-Refresh-Token-Speicherung.
  \item \textbf{Progressive Web App (PWA):} Durch ein Service Worker-Manifest
    und eine \texttt{manifest.webmanifest}-Datei kann die EHAE Web App als
    PWA installiert werden -- ohne App Store, direkt aus dem Intranet-Browser.
    Besonders relevant für Windows-Arbeitsstationen im Büro/Engineering-Umfeld,
    die keine dedizierte App-Installation erlauben.
  \item \textbf{Desktop-App via Electron/Tauri:} Für die Engineer App, die
    auf großen Monitoren mit Maus- und Tastaturinteraktion genutzt wird, bietet
    Tauri (Rust-basiert, $\sim$3\,MB Binärgröße) eine schlanke Alternative zu
    Electron. Die gleiche Web-App-Codebase wird als native Desktop-Applikation
    für Windows, macOS und Linux verpackt. Native OS-Integrationen (Systemtray,
    Dateiexport, serieller Port für direkten EtherCAT-Zugriff) können per
    Tauri-Plugin ergänzt werden.
  \item \textbf{Chromebook / Android-Enterprise-Deployment:} In
    Unternehmensumgebungen mit Google Workspace lassen sich EHAE-Apps
    über die Android-Enterprise-Management-API (AMAPI) zentral deployen,
    konfigurieren und in einem Kiosk-Modus sperren -- ohne manuelle
    Geräteeinrichtung durch den Anlagenbetreiber.
\end{itemize}

\subsubsection*{Standardisierung und ETG-Roadmap}

Der langfristig bedeutsamste Schritt für EHAE ist die formale Aufnahme in den
EtherCAT-Standard:

\begin{itemize}[leftmargin=2em]
  \item \textbf{ETG.1000-Erweiterung (Application Layer):} Die vorliegende
    Arbeit liefert alle notwendigen technischen Grundlagen für einen
    ETG-Proposal. Als nächster Schritt ist die Einreichung bei der ETG
    Technical Working Group sowie die Durchführung eines
    Multi-Vendor-Interoperabilitätstests (IOT) mit mindestens drei
    unabhängigen HAG-Implementierungen anzustreben.
  \item \textbf{Konformitätstest-Suite:} Analog zu den bestehenden
    ETG-Conformance-Tests für Slaves sollte eine EHAE-Conformance-Test-Suite
    entwickelt werden, die automatisiert prüft, ob ein HAG alle HCP-Nachrichten
    korrekt und vollständig implementiert. Dies ist Voraussetzung für ein
    offizielles ETG-Konformitätszertifikat.
  \item \textbf{Open-Source-Referenzimplementierung:} Die Veröffentlichung
    einer HAG-Refe\-renzimplementierung unter Apache-2.0-Lizenz (z.\,B.\ auf
    GitHub unter der ETG-Organisation) würde die Akzeptanz des Standards
    massiv erhöhen. Sie dient ETG-Mitgliedern als Integrationsgrundlage und
    der Forschungsgemeinschaft als reproduzierbares Bewertungssubstrat.
  \item \textbf{IEC-Normungsverfahren:} Mittelfristig sollte EHAE als
    Bestandteil von IEC~61784-5 (\emph{Installation profile for industrial
    networks over Ethernet}) eingereicht werden. Dies würde die
    Interoperabilität mit OPC~UA, PROFINET und EtherNet/IP auf normativer
    Ebene regeln und EHAE für Großunternehmen mit strikten
    Normierungsanforderungen (Automobilindustrie, Luftfahrt) öffnen.
\end{itemize}

\subsubsection*{Netzwerksicherheit und Zero Trust}

Das aktuelle Sicherheitsmodell basiert auf OAuth~2.0 mit PKCE und TLS~1.3 und
ist für den typischen Maschinenpark ausreichend. Für hochsensible Umgebungen
(kritische Infrastruktur, Rüstung, Pharma) sind folgende Erweiterungen zu
untersuchen:

\begin{itemize}[leftmargin=2em]
  \item \textbf{Zero-Trust-Netzwerkarchitektur:} Statt eines perimeterbasierten
    Sicherheitsmodells würde ein Zero-Trust-Ansatz (NIST SP~800-207) jeden
    HCP-Request unabhängig vom Netzwerkstandort des Clients authentifizieren
    und autorisieren. Dies erfordert eine kurzlebige, per-Request-Zertifikat-
    Ausstellung (SPIFFE/SPIRE) und eine kontinuierliche Anomalie-Überwachung
    des Datenverkehrs.
  \item \textbf{Post-Quanten-Kryptografie:} NIST hat 2024 die ersten
    Post-Quanten-Kryptografie-Standards (FIPS~203 ML-KEM, FIPS~204 ML-DSA)
    verabschiedet. Eine Erweiterung des HCP-Protokolls und der
    TLS-Konfiguration auf hybrid-Schlüsselaustausch (X25519 + ML-KEM-768)
    sollte mittelfristig erfolgen, um die Kommunikation gegenüber zukünftigen
    Quantencomputern zu schützen.
  \item \textbf{Hardware Security Module (HSM):} Für den HAG in
    sicherheitskritischen Anlagen sollte der private TLS-Schlüssel in einem
    dedizierten HSM (z.\,B.\ Infineon OPTIGA TPM, NXP SE050) gespeichert
    werden, das physische Tamper-Erkennung bietet und den Schlüssel auch
    bei Kompromittierung des Betriebssystems schützt.
  \item \textbf{Audit Log mit Blockchain-Anker:} Zur unveränderlichen
    Nachvollziehbarkeit aller Bedienereingriffe (wer hat wann welchen
    Sollwert gesetzt?) könnte ein Audit-Log-Modul Hashchains über
    alle HCP-Schreibkommandos bilden und regelmäßig auf einer permissioned
    Blockchain (Hyperledger Fabric) verankern -- relevant für
    regulierte Branchen mit GxP-Anforderungen (FDA 21 CFR Part 11).
\end{itemize}

\subsubsection*{Formale Verifikation und Zuverlässigkeitsanalyse}

Für den Einsatz in sicherheitsrelevanten Maschinen ist formale Verifikation
ein wichtiges Zukünftiges Arbeitsfeld:

\begin{itemize}[leftmargin=2em]
  \item \textbf{TLA\textsuperscript{+}-Modell des HCP-Protokolls:} Das
    HCP-Zustandsmaschine (Session-Aufbau, Token-Refresh, Alarm-Delivery) sollte
    in TLA\textsuperscript{+} modelliert und mit dem TLC Model Checker auf
    Deadlock- und Liveness-Eigenschaften geprüft werden. Dies liefert einen
    formalen Beweis, dass z.\,B.\ kein Zustand erreichbar ist, in dem ein
    Alarm dauerhaft nicht zugestellt wird.
  \item \textbf{Fuzzing und Property-based Testing:} HAG-Implementierungen
    sollten einem strukturierten Fuzzing-Programm (AFL++, libFuzzer) für
    den WebSocket-Parser und den JSON-Deserializer unterzogen werden, da
    malformierte Eingaben in Industrieumgebungen durch defekte Geräte
    realistisch auftreten. Ergänzend bietet sich Property-based Testing
    (Hypothesis, QuickCheck) für das Variable-Mapping-Subsystem an.
  \item \textbf{FMEA des Gesamtsystems:} Eine formale Failure Mode and Effects
    Analysis (FMEA) nach IEC~60812 für das EHAE-System würde die
    Schwachstellen-Priorisierung für zukünftige Härtungsmaßnahmen ermöglichen
    und ist Voraussetzung für eine Maschinenzulassung nach dem EU \emph{Machinery
    Regulation} 2023/1230.
\end{itemize}

% =====================================================================
%\section{Abkürzungsverzeichnis}
%
%\begin{acronym}
%\acro{EHAE}{EtherCAT HMI App Extension}
%\acro{HAG}{HMI Application Gateway}
%\acro{HCP}{HAG Communication Protocol}
%\acro{HMI}{Human-Machine Interface}
%\acro{PDO}{Process Data Object}
%\acro{SDO}{Service Data Object}
%\acro{RBAC}{Role-Based Access Control}
%\acro{FoE}{File over EtherCAT}
%\acro{TLS}{Transport Layer Security}
%\acro{JWT}{JSON Web Token}
%\acro{PKCE}{Proof Key for Code Exchange}
%\acro{TCO}{Total Cost of Ownership}
%\acro{VLAN}{Virtual Local Area Network}
%\acro{CoE}{CAN over EtherCAT}
%\acro{EoE}{Ethernet over EtherCAT}
%\acro{AoE}{ADS over EtherCAT}
%\end{acronym}

\begin{thebibliography}{99}
\bibitem{iec61158} IEC, \textit{IEC 61158: Industrial communication networks -- Fieldbus
  specifications}, 2019.
\bibitem{iec61784} IEC, \textit{IEC 61784-2: Additional fieldbus profiles for real-time
  networks}, 2019.
\bibitem{etg_spec} EtherCAT Technology Group, \textit{EtherCAT Specification: EAL},
  ETG.1000.6, 2022.
\bibitem{rfc6455} I.~Fette, A.~Melnikov, \textit{The WebSocket Protocol}, RFC~6455, IETF,
  2011.
\bibitem{rfc7636} N.~Sakimura et al., \textit{Proof Key for Code Exchange}, RFC~7636,
  IETF, 2015.
\bibitem{rfc8446} E.~Rescorla, \textit{TLS Protocol Version 1.3}, RFC~8446, IETF, 2018.
\bibitem{capacitor} Ionic, \textit{Capacitor: Cross-platform Native Runtime for Web Apps},
  \url{https://capacitorjs.com}, 2024.
\bibitem{flutter} Google LLC, \textit{Flutter: Build apps for any screen},
  \url{https://flutter.dev}, 2024.
\bibitem{maui} Microsoft, \textit{.NET MAUI Documentation},
  \url{https://learn.microsoft.com/dotnet/maui}, 2024.
\bibitem{wifi6} C.~Deng et al., ``IEEE 802.11ax,'' \textit{IEEE Communications Magazine},
  vol.~55, no.~12, 2017.
\bibitem{opc_ua} OPC Foundation, \textit{OPC UA Specification}, OPC 10000-1, 2022.
\bibitem{iso13849} ISO, \textit{ISO 13849-1: Safety of machinery -- Safety-related parts
  of control systems}, 2023.
\bibitem{etg_coe} EtherCAT Technology Group, \textit{CANopen over EtherCAT (CoE)},
  ETG.1000.6, 2022.
\end{thebibliography}

\end{document}
